\section{Self-Reference and Limits}\label{sec:self-reference-limits}

The Self-Reference and Limits chapter collects the laws that mark where
closure formation reaches a boundary.  Some boundaries are horizons:
boundary data can determine a declared target even when the ambient
carrier is larger.  Others concern information loss, singular limits,
locality, topology and completeness: what is visible, what is glued
globally, what survives a limit, and what cannot be established from
inside the same closure without additional audit structure.  The
self-audit law, Reflexive Nonclosure (\Fref{F21}), is stated in the
Status, Records, Coherence chapter because its data are claim records;
here the law that treats self-reference directly is the Closure
Completeness Boundary (\Fref{F52}).  In the running example of
\Cref{sec:intro}, the identification of \(1\) and \(2\) made by \(q\)
loses information that \(F\) records; in the language of the
Information Loss Normal Form (\Fref{F34}), with the identity as
transport, \(q\) as recovery quotient and \(F\) as source information,
\((1,2)\) lies in the loss obstruction.

The eighteen laws proceed from local sufficiency to completeness.
Horizon Sufficiency / Holographic Compression (\Fref{F33}) opens the
chapter: a target depends on a history only through its boundary
readout and exterior class exactly when it descends through their
pairing.  It is followed by Information Loss, Scale Descent /
Renormalization, Singular Limit Completion, Complementarity as
Non-Joint Access, Arrow from Record Monotonicity and Entropy as Fiber
Volume.  The middle of the chapter treats Irreducibility / No Short
Closed Description, Continuum Emergence, Criticality / Phase Boundary,
Locality / Domain of Dependence, Law--State Split and Fine-Tuning as
Small Selector Region.  The last part gives Topology as Global Gluing
Invariant, Common-Source / Nonlocal Correlation, Vacuum / Background
Selection, Unification as Common Refinement and the Closure
Completeness Boundary, which is distinct from the closure--reality
boundary of \Cref{sec:anchor:closure-reality-boundary}.  The chapter
depends only on \Cref{sec:apparatus}, and the summary table at its end
groups the laws into limits, emergence and invariants.

\subsection{Horizon Sufficiency / Holographic Compression [\texorpdfstring{\Fref{F33}}{F33}]}\label{sec:self-reference-limits:horizon-sufficiency}

\begin{definition}[Horizon-bounded support]\label{def:horizon-bounded-support}
Let \(\beta:I\to B\) be a boundary readout and \(\varepsilon:E\to Q\) an
exterior classifier on an interior--exterior decomposition \(I\times E\). A
target readout \(t:I\times E\to T\) is \emph{horizon-bounded} for
\((\beta,\varepsilon)\) when \(t(i,e)=t(i',e')\) whenever \(\beta(i)=\beta(i')\)
and \(\varepsilon(e)=\varepsilon(e')\); the pair \((\beta,\varepsilon)\) is its
horizon record.
\end{definition}

\begin{theorem}[Horizon Sufficiency / Holographic Compression]\label{thm:F33}
Let \(I\) and \(E\) be sets, the interior and exterior factors (a set \(H\) of histories with a decomposition \(H\cong I\times E\) is identified with \(I\times E\)), and define
\[
  \pi : I \times E \longrightarrow B \times Q,
  \qquad
  \pi(i,e)=(\beta(i),\varepsilon(e)),
\]
where \(\beta:I\to B\) is the boundary readout and
\(\varepsilon:E\to Q\) is the exterior classifier. Assume \(\pi\) is
surjective, and let \(t:I\times E\to T\) be a declared target readout.
Then
\begin{align*}
  t\ \text{horizon-bounded for }(\beta,\varepsilon)
  &\iff
  t \text{ descends through }\pi\\
  &\iff
  \forall (i,e),(i',e'),\
    \pi(i,e)=\pi(i',e')\Rightarrow t(i,e)=t(i',e')\\
  &\iff
  \exists\,\bar t:B\times Q\to T\ \text{ such that }\
    \bar t\circ\pi=t .
\end{align*}
Here horizon-boundedness is that of \Cref{def:horizon-bounded-support}. Call \(t\) \emph{horizon-sufficient}
when these hold. For a family \((t_\sigma)_{\sigma\in\Sigma}\), the bundled
readout \(x\mapsto(t_\sigma(x))_{\sigma}\) is horizon-sufficient if and only if
every \(t_\sigma\) is.
\end{theorem}
\begin{proof}
By \Cref{def:horizon-bounded-support}, \(t\) is horizon-bounded exactly
when it is constant on the fibers of \(\pi\).  Since \(\pi\) is
surjective, the descent criterion of \Cref{sec:apparatus} gives the
equivalence with descent through \(\pi\) and with the existence of
\(\bar t\).  Two points have the same bundled value exactly when they
have the same value under every \(t_\sigma\), so the bundled readout is
constant on the fibers of \(\pi\) exactly when every \(t_\sigma\) is.
\end{proof}

The motivating substrate-physics analogy is the holographic principle, in which
the bulk information of a region is encoded on its boundary
\citep{Maldacena1997LargeN,Bousso2002Holographic}.  The
layer-agnostic statement here is the descent-through-boundary form
of that pattern; it does not assert the AdS/CFT correspondence or
any particular gravitational realization.

\subsection*{Layer instantiations}\label{layer-inst:F33}

\begingroup\small
\begin{description}
\item[\emph{Boundary-value problem for a linear elliptic
{PDE} (mathematical analysis).}] \textit{Special case.}
For a fixed bounded Lipschitz domain \(\Omega\subseteq\mathbb R^n\), the
first factor \(I\) of the decomposition (its name in \Cref{thm:F33} is only a slot name: here the first factor carries the boundary data and the second factor carries the bulk data \(f\) and \(A\)) is the set of admissible Dirichlet data \(g\); the boundary readout records \(g\) in a fixed
function space (e.g., \(H^{1/2}(\partial\Omega)\)).
The second factor \(E\) parameterises sources \(f\in H^{-1}(\Omega)\) and bounded measurable, uniformly elliptic
divergence-form coefficient fields; the exterior classifier records the
source-and-coefficient data in their function-space
equivalence class.  The target readout is the entire
weak solution \(u\in H^1(\Omega)\).  By the Lax-Milgram
theorem, applied, after lifting \(g\) to \(H^1(\Omega)\), to the bilinear form
\(\int_\Omega A\nabla u\cdot\nabla v\), coercive on \(H^1_0(\Omega)\) by the Poincar\'e inequality, the target descends through the product
quotient: same boundary data plus same
source-coefficient class yield the same solution.  For Neumann data over a connected \(\Omega\), put
\(g\in H^{-1/2}(\partial\Omega)\) in the first factor \(I\) and
\(f\in(H^1(\Omega))'\), with the coefficient field, in the second factor \(E\),
and let the target be the weak solution modulo constants when
\(\langle f,1\rangle+\langle g,1\rangle_{\partial\Omega}=0\) and a marker
\(\bot\) otherwise. The compatibility test and, on compatible data,
coercivity on \(H^1(\Omega)/\mathbb R\) make this target a function of
\((g,f,A)\), so it again descends through the product quotient
\citep{Evans2010PartialDifferential}.

\item[\emph{Boundary representation in solid {CAD}
modelling (engineering).}] \textit{Analogy.}
A 3D rigid solid \(B\subset\mathbb R^3\) is represented
by its boundary surface together with adjacency and
orientation data: a B-rep records faces, edges, and
vertices with incidence relations.  The interior is the
open solid; the boundary readout records the B-rep with
face-orientation conventions; the exterior classifier
records the combinatorial topology (genus, handles,
through-holes).  Geometric and topological readouts ---
volume via the divergence theorem, mass distribution,
collision detection, boolean operations --- descend
through the product quotient
\citep{Mortenson2006GeometricModeling}.

\end{description}
\endgroup


\subsection{Information Loss Normal Form [\texorpdfstring{\Fref{F34}}{F34}]}\label{sec:self-reference-limits:information-loss}

\paragraph{Reading the statement.} The symbols are those of \Cref{thm:F34} below.
\(\operatorname{Recoverable}(\sigma,\tau,\rho)\) says that the source
information can be read back after transport: \(\sigma(h)\) depends only
on the later recovery class \(\rho(\tau h)\), that is, \(\sigma\)
descends through \(\rho\circ\tau\).

\begin{theorem}[Information Loss Normal Form]\label{thm:F34}
Let \(H_0,H_1\) be source and later carriers with transport
\(\tau:H_0\to H_1\). Let \(\sigma:H_0\to I\) be source information,
and \(\rho:H_1\to R\) a recovery quotient. Assume that
\(\rho\circ\tau\) is surjective onto \(R\). Define
\[
  \mathcal O_\rho
  =\{(h,h'):\rho(\tau h)=\rho(\tau h')\ \text{and}\
      \sigma(h)\ne\sigma(h')\}.
\]
Then
\begin{align*}
  \operatorname{Recoverable}(\sigma,\tau,\rho)
  &\iff \mathcal O_\rho=\emptyset\\
  &\iff
  \exists\,\bar\sigma:R\to I\ \text{ such that }\
    \bar\sigma\circ\rho\circ\tau=\sigma .
\end{align*}
The later quotient preserves the declared source information exactly
when every identification made by \(\rho\circ\tau\) is invisible to
\(\sigma\).
\end{theorem}
\begin{proof}
\(\operatorname{Recoverable}(\sigma,\tau,\rho)\) says that \(\sigma\)
descends through \(\rho\circ\tau\), and \(\mathcal O_\rho\) is the
split-pair obstruction of \(\rho\circ\tau\) and \(\sigma\).  Since
\(\rho\circ\tau\) is surjective onto \(R\), both equivalences are the
descent criterion of \Cref{sec:apparatus}.  A pair in
\(\mathcal O_\rho\) has the same recovered class but different source
information, which is precisely loss across a recovery class.
\end{proof}

A physical analogue is the black-hole information paradox, posed by Hawking \citep{Hawking1976Breakdown} after his thermodynamic analysis of black holes \citep{Hawking1976BHThermodynamics}, where the loss of information is argued
from semiclassical gravity; the statement above records only the descent
obstruction of a declared horizon quotient.

\subsection*{Layer instantiations}\label{layer-inst:F34}

\begingroup\small
\begin{description}
\item[\emph{Huffman entropy coding (information theory).}] \textit{Special case.}
The source carrier is the set of source-symbol sequences
over an alphabet \(\Sigma\); the source information is
the original symbol sequence.  The transport encodes
into a prefix-free binary codeword sequence; the later carrier is the set \(\tau(H_0)\) of encoded streams, and the recovery quotient is the prefix decoder on it.  Distinct
source sequences produce distinct codeword streams; the
obstruction set is empty and the source is fully
recoverable.  Compression achieves the entropy bound to
within one bit per symbol
\citep{Huffman1952Method}.

\item[\emph{Reed-Solomon forward error correction
(coding theory).}] \textit{Special case.}
For a Reed-Solomon \((n,k)\) code over a finite field
\(\mathbb F_q\), the source carrier is the set of pairs
\((m,\varepsilon)\) of a \(k\)-symbol message and an error
vector of Hamming weight at most
\(t=\lfloor(n-k)/2\rfloor\), with \(\sigma(m,\varepsilon)=m\).
Let the later carrier be the union of radius-\(t\) Hamming
balls around the codewords. The transport sends
\((m,\varepsilon)\) to \(c(m)+\varepsilon\), where \(c(m)\)
evaluates the message polynomial at \(n\) distinct field
points. Let \(\rho\) be bounded-distance decoding from this
carrier to the message set, using syndrome decoding and error
correction. Disjoint radius-\(t\) balls make
\(\rho(c(m)+\varepsilon)=m\), so the obstruction is empty;
taking \(\varepsilon=0\) for each message shows that
\(\rho\circ\tau\) is surjective, and the source message is
recoverable; Reed-Solomon codes underlie CD/DVD, QR
codes, and deep-space communication
\citep{ReedSolomon1960Codes}.

\item[\emph{Cryptographic one-way hash function
(cryptography).}] \textit{Special case.}
For a cryptographic hash function
\(H:\{0,1\}^{<2^{64}}\to\{0,1\}^{256}\) (SHA-256), the
source information is the original bitstring.  The
transport maps bitstrings to fixed-length digests;
the recovery quotient is the identity on the set \(H(\{0,1\}^{<2^{64}})\) of attained digests.
Since the domain has more than \(2^{256}\) elements, distinct source
bitstrings collide to the same digest by pigeonhole, so the obstruction set is non-empty.  The
cryptographic hash is the canonical non-recoverable
instance; its cryptographic use rests on collisions
being hard to find, not absent
\citep{NIST2015FIPS180}.

\item[\emph{Database write-ahead logging and {ARIES}
recovery (database systems).}] \textit{Analogy.}
The source information is the committed-state effects
of database transactions.  The transport writes log
records (before-image, after-image, LSN, transaction
ID) to stable storage before any modified page is
written.  The recovery quotient is the ARIES
three-pass algorithm (analysis, redo, undo) applied
to the on-disk log and data pages after a crash.
Within the failure model, the obstruction set is
empty and the committed-state effects are recoverable
\citep{Mohan1992ARIES}.

\item[\emph{Magnetic-resonance image reconstruction
(medical imaging).}] \textit{Analogy.}
The source information is the tissue proton-density
map, supported within the field of view and represented at the
declared resolution.  The transport
samples k-space (the Fourier dual of physical space)
under controlled magnetic-field gradient sequences;
the later carrier is the set of sampled k-space
datasets.  The recovery quotient is the inverse
Fourier transform, in clinical practice combined with
parallel-imaging unfolding and partial-Fourier
filling.  When the k-space sampling step is at most the reciprocal of the
object's spatial extent, the obstruction set is empty and the
proton-density map is recoverable up to the resolution set by the
sampled k-space extent; with a coarser step, aliasing folds the
object onto itself and makes the loss illegitimate
\citep{Lauterbur1973MRI}.
\end{description}
\endgroup

\subsection{Scale Descent / Renormalization [\texorpdfstring{\Fref{F35}}{F35}]}\label{sec:self-reference-limits:scale-descent}

\begin{theorem}[Scale Descent / Renormalization]\label{thm:F35}
Let \(H\) be a carrier with consecutive scale quotients
\(q_n:H\to Q_n\) and \(q_{n+1}:H\to Q_{n+1}\), with \(q_n\) surjective. Let
\(c:Q_n\to Q_{n+1}\) be a surjective coarsening with
\[
  c\circ q_n=q_{n+1},
\]
and let \(t:H\to T\) be a target law. Put
\[
  \mathcal O_n(t)=
  \{(h,h'):q_n(h)=q_n(h')\ \text{and}\ t(h)\ne t(h')\}.
\]
Then
\begin{align*}
  t \text{ descends to } Q_n
  &\iff \mathcal O_n(t)=\emptyset\\
  &\iff
  \exists\,\bar t_n:Q_n\to T\ \text{ such that }\
    \bar t_n\circ q_n=t .
\end{align*}
When \(t\) descends to \(Q_n\), the map \(\bar t_n\) is unique, and a
renormalized law at scale \(Q_{n+1}\), a map \(\bar t_{n+1}:Q_{n+1}\to T\) with
\(\bar t_{n+1}\circ c=\bar t_n\), exists exactly when \(\bar t_n\) is constant
on the fibers of \(c\), and then \(t=\bar t_{n+1}\circ q_{n+1}\), and \(\bar t_{n+1}\) is unique since \(c\) is surjective. Consequently
\(t\) descends to \(Q_{n+1}\) if and only if it descends to \(Q_n\) and
\(\bar t_n\) is constant on the fibers of \(c\).
\end{theorem}
\begin{proof}
Since \(q_n\) is surjective, the first two equivalences and the
uniqueness of \(\bar t_n\) are the descent criterion of
\Cref{sec:apparatus} for \(q_n\) and \(t\). For the next scale, if \(\bar t_{n+1}\circ c=\bar t_n\),
then \(c(x)=c(x')\) gives \(\bar t_n(x)=\bar t_n(x')\), so \(\bar t_n\) is
constant on the fibers of \(c\). Conversely, if \(\bar t_n\) is constant on
the fibers of \(c\), then \(\bar t_{n+1}(c(x)):=\bar t_n(x)\) is independent of
the choice of \(x\), and surjectivity of \(c\) defines it on all of
\(Q_{n+1}\). Then \(t=\bar t_n\circ q_n=\bar t_{n+1}\circ c\circ q_n=\bar t_{n+1}\circ q_{n+1}\).
Conversely, if \(t=s\circ q_{n+1}=s\circ c\circ q_n\), then \(t\) descends to
\(Q_n\), uniqueness gives \(\bar t_n=s\circ c\), and \(\bar t_n\) is constant on
the fibers of \(c\). If two renormalized laws \(\bar t_{n+1},\bar t'_{n+1}\)
satisfy \(\bar t_{n+1}\circ c=\bar t_n=\bar t'_{n+1}\circ c\), then they agree on
the image of \(c\), which is all of \(Q_{n+1}\). Thus scale descent is quotient
descent iterated along the tower.
\end{proof}

The Wilsonian renormalization group of statistical and quantum field
theory \citep{Wilson1975RG} is the standard substrate analogy, in
which scale-by-scale quotient descent is the central calculational
device.

\subsection*{Layer instantiations}\label{layer-inst:F35}

\begingroup\small
\begin{description}
\item[\emph{Hydrodynamic limit via Chapman-Enskog
expansion (kinetic theory).}] \textit{Analogy.}
For a rarefied gas, the carrier is the space of
phase-space distribution functions \(f(x,v,t)\).  The
scale quotient at the kinetic scale records the
one-particle distribution at Knudsen number
\(\mathrm{Kn}\sim O(1)\); the coarsening sends \(f\)
to its hydrodynamic moments (density, momentum
density, energy density) in the limit
\(\mathrm{Kn}\to 0\).  The target law at the kinetic
scale is the Boltzmann equation; the renormalised law
at the hydrodynamic scale is the Navier-Stokes (or
Euler) system, derived by the Chapman-Enskog
asymptotic expansion.  Scale descent succeeds in
\(\mathrm{Kn}\ll 1\) regimes and fails near shocks
and boundary layers \citep{Cercignani1988Boltzmann}.

\item[\emph{Burt-Adelson Laplacian image pyramid
(computer vision and signal processing).}] \textit{Special case.}
For a digital image, the carrier is the space of
full-resolution images.  The scale quotient
sends an image to its Gaussian-pyramid level at
resolution \(2^{-n}\) (Gaussian filter followed by
downsampling), corestricted to the levels it attains; the coarsening sends level \(n\) to
level \(n+1\) by repeating the filter-and-downsample
step.  The band-pass Laplacian layer
\(L_n=G_n-\mathrm{up}(G_{n+1})\) is the target law at
level \(n\) and descends through the scale quotient.
The Laplacian pyramid is a complete representation:
the original image \(G_0\) is recovered exactly from the coarsest level \(G_N\) and the residues \(L_0,\dots,L_{N-1}\) by \(G_n=L_n+\mathrm{up}(G_{n+1})\), \(n=N-1,\dots,0\)
\citep{BurtAdelson1983LaplacianPyramid}.

\end{description}
\endgroup


\subsection{Singular Limit Completion [\texorpdfstring{\Fref{F36}}{F36}]}\label{sec:self-reference-limits:singular-limit}

\paragraph{Reading the statement.} In \Cref{thm:F36} below, a
limiting process \(p\) has a completed object \(\ell(p)\), and
\(\operatorname{Completed}(u,\ell)\) says that its target-limit value
depends only on that object.  \(\operatorname{PathDependent}(u,\ell)\)
says that two processes with the same completed object have different
target-limit values.  \(\operatorname{Removable}(j,\ell,\sigma,p_0)\)
says that completion removes the singularity at \(p_0\): the completed
object \(\ell(p_0)\) is not the image of a current object under \(j\),
and the signature-limit readout \(\sigma\) descends through \(\ell\).

\begin{theorem}[Singular Limit Completion]\label{thm:F36}
Let \(P\) be a set of limiting processes, let \(\ell:P\to\widehat Q\)
be a surjective limit quotient, let \(j:Q\to\widehat Q\) be a
completion map, let \(u:P\to\widehat T\) be the target-limit readout,
let \(\sigma:P\to\widehat S\) be a declared signature-limit readout,
let \(u^{\mathrm{ren}}:P\to\widehat T^{\mathrm{ren}}\) be a declared
renormalized target-limit readout, and let \(p_0\in P\) be a selected
process.  Define the singular target obstruction
\[
  \mathcal O_{\ell}(u)
  =\{(p_1,p_2)\in P^2:
      \ell(p_1)=\ell(p_2)\ \text{and}\ u(p_1)\ne u(p_2)\}
\]
and the corresponding signature obstruction
\(\mathcal O_{\ell}(\sigma)\).  Write
\[
  \rho_{u,\ell}:P\to\widehat Q\times\widehat T,
  \qquad
  \rho_{u,\ell}(p):=(\ell(p),u(p)),
\]
for the joint \((\ell,u)\) refinement, and let
\(\operatorname{Current}_{j}(\ell,p_0)\) abbreviate
\(\exists c\in Q,\ j(c)=\ell(p_0)\), and define
\[
  \operatorname{Removable}(j,\ell,\sigma,p_0)
  :\Longleftrightarrow
  \neg\operatorname{Current}_{j}(\ell,p_0)
  \wedge
  \operatorname{Completed}(\sigma,\ell).
\]
Here \(\operatorname{Completed}(v,\kappa)\), for a readout \(v\) and a map
\(\kappa\) on \(P\), says that \(v(p)\) depends only on \(\kappa(p)\),
\(\operatorname{Refines}(\kappa,\ell)\) says that \(\kappa(p_1)=\kappa(p_2)\)
implies \(\ell(p_1)=\ell(p_2)\), and
\(\operatorname{PathDependent}(u,\ell)\) says that
\(\mathcal O_\ell(u)\ne\emptyset\).
Then
\begin{align*}
  \operatorname{Completed}(u,\ell)
  &\iff \mathcal O_{\ell}(u)=\emptyset,\\
  \operatorname{Completed}(u,\ell)
  &\iff
  \exists\,\bar u:\widehat Q\to\widehat T,\ \bar u\circ\ell=u,\\
  \operatorname{Completed}(\sigma,\ell)
  &\iff \mathcal O_{\ell}(\sigma)=\emptyset,\\
  \operatorname{Removable}(j,\ell,\sigma,p_0)
  &\iff \neg\operatorname{Current}_{j}(\ell,p_0)\wedge\mathcal O_{\ell}(\sigma)=\emptyset,\\
  \operatorname{PathDependent}(u,\ell)
  &\iff \neg\operatorname{Completed}(u,\ell),\\
  &\operatorname{Refines}(\rho_{u,\ell},\ell),\\
  &\operatorname{Completed}(u,\rho_{u,\ell}),\\
  \operatorname{Completed}(u^{\mathrm{ren}},\ell)
  &\iff
  \mathcal O_{\ell}(u^{\mathrm{ren}})=\emptyset.
\end{align*}
Equivalently, completion is descent of \(u\) through \(\ell\); the
canonical path-memory repair \(\rho_{u,\ell}\) is a refinement of
\(\ell\) that completes \(u\); a renormalized target is tested by the
same descent criterion as its raw target, and if \(u^{\mathrm{ren}}=R\circ u\) for a
map \(R\), then \(\operatorname{Completed}(u,\ell)\) implies
\(\operatorname{Completed}(u^{\mathrm{ren}},\ell)\), with the converse when \(R\) is
injective.
\end{theorem}
\begin{proof}
We prove the eight conjuncts. Parts (4) and (5) prove the PathDependent and Removable lines, which the statement lists in the opposite order; part (8) also proves the clause on \(u^{\mathrm{ren}}=R\circ u\).

\emph{(1)--(3) Completion and obstructions.}
Since \(\ell\) is surjective, the descent criterion of
\Cref{sec:apparatus}, applied to \(\ell\) and \(u\), gives (1) and (2),
and applied to \(\ell\) and \(\sigma\) gives (3).

\emph{(4) Path-dependent iff not completed.}
Path dependence is witnessed by a split pair:
\[
  \operatorname{PathDependent}(u,\ell)
  \iff
  \exists(p_1,p_2)\in\mathcal O_{\ell}(u).
\]
By (1), \(\mathcal O_{\ell}(u)\ne\emptyset\) is equivalent to
\(\neg\operatorname{Completed}(u,\ell)\).

\emph{(5) Removable.}
By definition \(\operatorname{Removable}(j,\ell,\sigma,p_0)\) is
\(\neg\operatorname{Current}_{j}(\ell,p_0)\wedge\operatorname{Completed}(\sigma,\ell)\),
and (3) replaces the second conjunct by \(\mathcal O_{\ell}(\sigma)=\emptyset\).

\emph{(6) The path repair refines \(\ell\).}
The canonical path-memory repair is
\[
  \rho_{u,\ell}(p)=(\ell(p),u(p)).
\]
Projection to the first coordinate recovers the original limit:
\[
  \operatorname{pr}_{\widehat Q}\circ\rho_{u,\ell}=\ell.
\]
Thus each \(\rho_{u,\ell}\)-fiber lies inside an \(\ell\)-fiber, so
\(\rho_{u,\ell}\) is a source refinement of \(\ell\).

\emph{(7) After the path repair, the target is completed.}
Projection to the second coordinate recovers the target-limit readout:
\[
  \operatorname{pr}_{\widehat T}\circ\rho_{u,\ell}=u.
\]
Hence \(u\) factors through the repaired quotient by the explicit
factor \(\operatorname{pr}_{\widehat T}\), and therefore
\[
  \operatorname{Completed}(u,\rho_{u,\ell}).
\]

\emph{(8) Renormalized completed iff its obstruction is empty.}
The descent criterion applied to \(\ell\) and \(u^{\mathrm{ren}}\) gives
\[
  \operatorname{Completed}(u^{\mathrm{ren}},\ell)
  \iff
  \mathcal O_{\ell}(u^{\mathrm{ren}})=\emptyset,
\]
so a renormalized target is subject to the same completion criterion
as a raw target. If \(u^{\mathrm{ren}}=R\circ u\) and \(u=\bar u\circ\ell\), then
\(u^{\mathrm{ren}}=(R\circ\bar u)\circ\ell\); conversely, assume \(R\) is injective
and \(\operatorname{Completed}(u^{\mathrm{ren}},\ell)\). For \(\ell(p_1)=\ell(p_2)\),
this completion gives \(R(u(p_1))=R(u(p_2))\); injectivity gives
\(u(p_1)=u(p_2)\). Thus \(\mathcal O_\ell(u)=\emptyset\), and the first
equivalence gives \(\operatorname{Completed}(u,\ell)\).
\end{proof}

\subsection*{Layer instantiations}\label{layer-inst:F36}

\begingroup\small
\begin{description}
\item[\emph{Riemann removable singularity in complex
analysis (complex analysis).}] \textit{Special case.}
For a function \(f\) holomorphic on a punctured disk,
the set of limiting processes is the family of
approach paths to the puncture along which \(f\) has a
limit in the Riemann sphere; the limit quotient
\(\ell\) is constant, since every path approaches the same puncture; the
target readout is the path-limit value of \(f\).  By
Riemann's removable-singularity theorem, if \(f\) is
bounded on the punctured disk then the limit is
independent of approach path and the obstruction set
is empty.  At a pole
\(|f|\to\infty\) along every approach path, so the limit is
path-independent on the Riemann sphere.  An essential
singularity can give the path-dependent case: \(e^{1/z}\) tends
to \(\infty\) along the positive real axis and to \(0\) along the
negative real axis
\citep{Ahlfors1979ComplexAnalysis}.

\item[\emph{L'H\^opital's rule for indeterminate
forms (real analysis).}] \textit{Analogy.}
For functions \(f,g\) with \(f(a)=g(a)=0\), the
literal substitution \(\lim_{x\to a}f(x)/g(x)\)
yields the indeterminate form \(0/0\), which by
itself does not descend through the limit quotient.
L'H\^opital's rule provides the completion: if \(f\)
and \(g\) are differentiable with \(g'(a)\ne 0\),
the limit equals \(f'(a)/g'(a)\).  The path-memory
repair records the order of vanishing of \(f\) and
\(g\) at \(a\), completing the limit through the
derivative-ratio readout
\citep{Cauchy1821CoursAnalyse}.

\item[\emph{Riemann integral as limit of Riemann sums
(classical analysis).}] \textit{Special case.}
For a bounded \(f:[a,b]\to\mathbb R\), the set of
limiting processes is the family of sequences of
tagged partitions with mesh tending to zero along
which the Riemann sums converge; the target readout
is the limit of the Riemann sums; the limit quotient \(\ell\) is the constant
map onto a one-point set.  The function is Riemann-integrable precisely when
\(\mathcal O_\ell(u)=\emptyset\) (for bounded \(f\), tags approximating the interval suprema
and infima give sequences converging to the upper and lower
Darboux integrals, which differ otherwise), i.e., when the limit of Riemann sums
is independent of the choice of partition
refinement and tag selection.  The completed map is
the Riemann integral
\(\int_a^b f\,dx\); non-Riemann-integrable functions
(e.g., the Dirichlet function) instantiate the
path-dependent case
\citep{Riemann1854Habilitationsschrift}.
\end{description}
\endgroup


\subsection{Complementarity as Non-Joint Access [\texorpdfstring{\Fref{F37}}{F37}]}\label{sec:self-reference-limits:complementarity}

Position and momentum \citep{Bohr1928Quantum} motivate a restricted choice
of \(\mathcal J\): if an admissible carrier must arise from a single joint
quantum observable with sharp position and sharp momentum as its marginals, no
such carrier exists. The pair of statistical readouts still exists as a map.
Merely requiring a joint probability distribution with the correct marginals
would not give non-joint access, since the product of those marginals is one.
The statement below abstracts non-joint access relative to the declared class
\(\mathcal J\).

\paragraph{Reading the statement.}
The class \(\mathcal J\) and the formedness of each access are declared
data: they fix which joint carriers count as admissible at the stated
scope. \(q_A\perp_{\mathrm{acc}}q_B\) says that each access is lawful
on its own but no admissible joint carrier determines both.

\begin{theorem}[Complementarity as Non-Joint Access]\label{thm:F37}
Let \(H\) be a nonempty carrier with two access quotients
\(q_A:H\to Q_A\) and \(q_B:H\to Q_B\) and target readouts
\(t_A:H\to T_A\), \(t_B:H\to T_B\).  The access \(q_A\) is
\emph{separately admissible} when \(t_A\) descends through \(q_A\)
(formedness of the access is part of the setting and adds no condition), and
likewise for \(q_B\).  Let \(K\) be a declared
joint-carrier type and let \(\mathcal J\) be a declared class of
admissible joint carriers \(j:H\to K\).  A joint access in
\(\mathcal J\) is a carrier \(j\in\mathcal J\) together with maps
\(a:K\to Q_A\) and \(b:K\to Q_B\) satisfying
\[
  a\circ j=q_A,\qquad b\circ j=q_B .
\]
Define
\begin{align*}
  q_A \perp_{\mathrm{acc}} q_B
  \;:\Longleftrightarrow\;&
  q_A,\,q_B\ \text{separately admissible}\\
  &\wedge\ \neg\exists\, j\in\mathcal J,\ \exists\,a:K\to Q_A,\ \exists\,b:K\to Q_B,\
  a\circ j=q_A\ \wedge\ b\circ j=q_B .
\end{align*}
Then \(q_A \perp_{\mathrm{acc}} q_B\) holds if and only if \(q_A\) and \(q_B\) are separately admissible and every
\(j\in\mathcal J\) identifies two histories that one of the accesses
separates: for each \(j\in\mathcal J\) there are \(h,h'\in H\) with
\(j(h)=j(h')\) and (\(q_A(h)\neq q_A(h')\) or \(q_B(h)\neq q_B(h')\)).
Thus complementarity is precisely separate admissibility together with the presence, in every admissible joint
carrier, of a split pair against one of the two accesses. Equivalently, \(q_A\perp_{\mathrm{acc}}q_B\) holds if and only if \(q_A\) and \(q_B\) are separately admissible and no \(j\in\mathcal J\) has
\(\langle q_A,q_B\rangle=g\circ j\) for a map \(g:K\to Q_A\times Q_B\); in
particular, if \(K=Q_A\times Q_B\) and \(\langle q_A,q_B\rangle\in\mathcal J\),
then \(q_A\perp_{\mathrm{acc}}q_B\) fails.
\end{theorem}
\begin{proof}
Fix
\(j\in\mathcal J\). If \(a\circ j=q_A\), then \(j(h)=j(h')\) gives
\(q_A(h)=q_A(h')\). (For the factorization form, take \(g=\langle a,b\rangle\), and
conversely \(a=\mathrm{pr}_1\circ g\), \(b=\mathrm{pr}_2\circ g\); if the pairing
itself lies in \(\mathcal J\), take \(g=\mathrm{id}\).) Conversely, if \(q_A\) is constant on the fibers of
\(j\), set \(a(j(h)):=q_A(h)\) on the image of \(j\) and give \(a\) any value
in \(Q_A\) elsewhere, which is possible since \(H\), and hence \(Q_A\), is
nonempty; then \(a\circ j=q_A\). So a map \(a\) exists exactly when no pair
identified by \(j\) is separated by \(q_A\), and likewise for \(b\) and
\(q_B\). Hence \(j\) is a joint access exactly when it identifies no pair that
\(q_A\) or \(q_B\) separates, and \(\mathcal J\) contains no joint access
exactly when every \(j\in\mathcal J\) identifies such a pair. Together with
separate admissibility this is complementarity. The product map
\(h\mapsto(q_A(h),q_B(h))\) always carries both accesses, and it factors
through every carrier of both accesses; complementarity is
therefore a statement about the declared class \(\mathcal J\), which
decides whether such a carrier counts as an admissible joint access.
\end{proof}

\subsection*{Layer instantiations}\label{layer-inst:F37}

\begingroup\small
\begin{description}
\item[\emph{Time-frequency Gabor uncertainty (signal
processing).}] \textit{Analogy.}
For a square-integrable signal \(s\in L^2(\mathbb R)\),
the access quotient \(q_A\) records its time-marginal
\(|s(t)|^2\) and \(q_B\) records its
frequency-marginal \(|\hat s(f)|^2\).  Within the
admissibility class of Cohen-class bilinear
time-frequency representations satisfying the correct
marginals, the covariance properties of bilinear
distributions, and non-negativity (probability-density
interpretation), no joint quotient exists: the
Wigner-Ville distribution satisfies marginals and
covariance but is not non-negative, and no member of
Cohen's class is jointly compatible.  The access
quotients are \Fref{F37}-complementary
\citep{Cohen1995TimeFrequency}.

\end{description}
\endgroup


\subsection{Arrow from Record Monotonicity [\texorpdfstring{\Fref{F38}}{F38}]}\label{sec:self-reference-limits:arrow}

The phrase ``arrow of time'' is due to \citet{Eddington1928Nature}, who tied it to the increase of entropy; the formulation below takes the arrow instead from monotone growth of records, without committing to a specific physical realization.

\paragraph{Reading the statement.} The symbols are those of \Cref{thm:F38} below.
\(\operatorname{Arrow}(\kappa)\) is defined at the level of histories:
the continuation respects records, and record status does not decrease along
it. The theorem says that this is the same as a map \(\bar\kappa\) on records
along which status does not decrease.

\begin{theorem}[Arrow from Record Monotonicity]\label{thm:F38}
Let \(H_s,H_t\) be source and target history sets with continuation
\(\kappa:H_s\to H_t\). Let \(r_s:H_s\to R_s\) be a surjective record quotient
and \(r_t:H_t\to R_t\) a target record map, and let
\(s_s:R_s\to S\), \(s_t:R_t\to S\) be status readouts into a set \(S\) with a relation \(\le_S\) (no order axioms are used). Define
\[
  \operatorname{Arrow}(\kappa)
  :\Longleftrightarrow
  r_t\circ\kappa\ \text{descends through}\ r_s
  \ \wedge\
  \forall h\in H_s,\ s_s(r_s h)\le_S s_t(r_t(\kappa h)).
\]
Then
\begin{align*}
  \bigl(\forall h,h'\in H_s,\ r_s(h)=r_s(h')\Rightarrow r_t(\kappa h)=r_t(\kappa h')\bigr)
  &\iff
  \exists\,\bar\kappa:R_s\to R_t\ \text{ such that }\
    \bar\kappa\circ r_s=r_t\circ\kappa,\\
  \operatorname{Arrow}(\kappa)
  &\iff
  \exists\,\bar\kappa:R_s\to R_t,\
  \bar\kappa\circ r_s=r_t\circ\kappa\\
  &\hphantom{{}\iff{}}\text{and } s_s(x)\le_S s_t(\bar\kappa x)\
    \forall x\in R_s.
\end{align*}
The arrow is the monotone direction of the descended record transport.
\end{theorem}
\begin{proof}
Since \(r_s\) is surjective, the first equivalence is the descent
criterion of \Cref{sec:apparatus} for \(r_s\) and \(r_t\circ\kappa\),
and the descended map \(\bar\kappa\) is unique.  The arrow condition
therefore splits into descent of \(r_t\circ\kappa\) and a status
inequality.  For \(x=r_s(h)\), \(s_t(\bar\kappa x)=s_t(r_t(\kappa h))\); since \(r_s\) is
surjective, the history-level and record-level inequalities are the same
condition, which gives the second equivalence.
\end{proof}

\subsection*{Layer instantiations}\label{layer-inst:F38}

\begingroup\small
\begin{description}
\item[\emph{Git commit {DAG} and version-control history
(software engineering).}] \textit{Analogy.}
The source history set is a project-state at a given
commit; the target is the project-state after a child
commit; \(\kappa\) records the commit operation.  The
record quotient sends a state to its commit hash, with
status readout equal to the DAG-ancestry generation
(the longest ancestor chain length).  The hash of a
new commit is the cryptographic digest of its content
together with parent-commit references, so the
descended map carries each parent hash to the child
hash.  The status readout is strictly monotone: each
child commit has generation greater than each of its
parents \citep{ChaconStraub2014ProGit}.

\end{description}
\endgroup


\subsection{Entropy as Fiber Volume [\texorpdfstring{\Fref{F39}}{F39}]}\label{sec:self-reference-limits:entropy}

The fiber-volume reading of entropy combines two classical substrate
sources: the statistical-mechanical interpretation
\(S=k\log W\) developed in \citet{Boltzmann1877Beziehung} and the
communication-theoretic entropy of \citet{Shannon1948Mathematical}.
The layer-agnostic statement below records the structural form
common to both.

\paragraph{Reading the statement.} In \Cref{thm:F39} below,
\(\operatorname{TargetEntropyZero}(q,t)\) is defined combinatorially,
as \(\mathcal O_q(t)=\emptyset\): no \(q\)-fiber contains two different
target values.  Its identification with a vanishing entropy value is
the added clause on \(m(\langle q,t\rangle,q)\).  The fiber-volume
functional of the last clause is the conditional Shannon entropy of a
uniformly distributed history given its class, and \(m(q_F,q_C)\) is
then the entropy of the \(q_F\)-class given the \(q_C\)-class, so
\(E(q_C)=E(q_F)+m(q_F,q_C)\) is the chain rule of conditional entropy;
the theorem evaluates \(m\) only on refining pairs.  In this instance
\(m(\langle q,t\rangle,q)\) is the conditional entropy of the target
given the \(q\)-class, and the last clause identifies
\(\operatorname{TargetEntropyZero}(q,t)\) with its vanishing.

\begin{theorem}[Entropy as Fiber Volume]\label{thm:F39}
Let \(q:H\to Q\) be a surjective access quotient with target readout
\(t:H\to T\).  Let \(q_F:H\to Q_F\) be a finer quotient and
\(q_C:H\to Q_C\) a coarsening, related by a factorization
\(c:Q_F\to Q_C\) with
\[
  c\circ q_F=q_C.
\]
Let \(\mathcal E\) be a set of entropy values with a relation \(\le_E\) (no order axioms are used), and let
\[
  E:\{\,q':H\to Q'\ :\ Q'\text{ a set of a fixed universe}\,\}\longrightarrow \mathcal E
\]
be an entropy functional assigning a fiber-volume value to every
declared quotient.  Suppose the apparatus declares a mixing
functional \(m\) and a combiner
\(\odot:\mathcal E\times\mathcal E\to\mathcal E\) satisfying the
chain-rule identity
\[
  E(q_C)=E(q_F)\odot m(q_F,q_C)
\]
and the combine-monotonicity condition at this pair
\[
  E(q_F)\le_E E(q_F)\odot m(q_F,q_C)
\]
(which holds when \(a\le_E a\odot b\) for all \(a,b\), and in the fiber-volume
instance because \(m\ge0\)).
Suppose also that \(\mathcal E\) has a declared zero \(0\) such that
\(E(q')=0\) exactly when \(q'\) is injective, that is, every nonempty fiber of \(q'\) is a singleton.
Define the raw split-pair and target obstructions
\begin{align*}
  \mathcal O_q^{\mathrm{raw}}
  &:=
  \{(h,h')\in H^2:q(h)=q(h')\ \text{and}\ h\ne h'\},\\
  \mathcal O_q(t)
  &:=
  \{(h,h')\in H^2:q(h)=q(h')\ \text{and}\ t(h)\ne t(h')\},
\end{align*}
and put \(\operatorname{TargetEntropyZero}(q,t):\Longleftrightarrow\mathcal O_q(t)=\emptyset\)
and \(\operatorname{TargetUnresolved}(q,t):\Longleftrightarrow\mathcal O_q(t)\neq\emptyset\);
the entropy identification is the clause on \(m(\langle q,t\rangle,q)\) below.
Of the displayed lines, only the first and the last involve \(E\); the target predicates involve \(E\) only through the mixing clause below.
Then
\begin{align*}
  E(q)=0
  &\iff
  \mathcal O_q^{\mathrm{raw}}=\emptyset,\\
  \operatorname{TargetEntropyZero}(q,t)
  &\iff
  t\ \text{descends through}\ q,\\
  \operatorname{TargetUnresolved}(q,t)
  &\iff
  \neg(t\ \text{descends through}\ q),\\
  \forall b\in Q_C,\ \forall h\in H,\quad
  q_C(h)=b
  &\iff
  \exists a\in Q_F,\ c(a)=b\ \wedge\ q_F(h)=a,\\
  E(q_F)
  &\le_E
  E(q_C).
\end{align*}
If moreover \(m(\langle q,t\rangle,q)=0\) exactly when \(\langle q,t\rangle\)
and \(q\) have the same fibers (as holds when \(m(q',q'')=0\) exactly when
\(q'\) and \(q''\) have the same fibers, for every \(q'\) refining \(q''\)), then
the target entropy is a mixing value:
\[
  m(\langle q,t\rangle,q)=0
  \iff t\ \text{descends through}\ q
  \iff \operatorname{TargetEntropyZero}(q,t).
\]
For finite nonempty \(H\), the fiber-volume functional
\(E(q')=|H|^{-1}\sum_{h\in H}\log|q'^{-1}(q'(h))|\) with
\(m(q',q'')=E(q'')-E(q')\) for \(q'\) refining \(q''\), with values in
\(\mathbb R\) under addition and zero \(0\), satisfies the zero criterion, and the chain-rule identity and
combine-monotonicity at every refining pair,
\(m\ge0\), and \(m(q',q'')=0\) exactly when \(q'\) and \(q''\) have the same
fibers; hence all conclusions above hold for it, in particular
\(E(q_F)\le E(q_C)\) and
\(\operatorname{TargetEntropyZero}(q,t)\iff m(\langle q,t\rangle,q)=0\).
\end{theorem}
\begin{proof}
Assume \(q\) is surjective, \(c\circ q_F=q_C\), the chain-rule
identity \(E(q_C)=E(q_F)\odot m(q_F,q_C)\), and combine-monotonicity
\(E(q_F)\le_E E(q_F)\odot m(q_F,q_C)\).  We prove the five conjuncts and then
the added clause.

\emph{(1) Raw entropy zero iff no raw split pair.}
By the declared zero, \(E(q)=0\) exactly when every \(q\)-fiber is a
singleton, that is, when no \(q\)-fiber contains distinct elements,
which is \(\mathcal O_q^{\mathrm{raw}}=\emptyset\).

\emph{(2) Target entropy zero iff target descends.}
Target entropy is measured by the target obstruction
\(\mathcal O_q(t)\).  Its emptiness is exactly the quotient-descent
condition for \(t\) through \(q\):
\[
  \mathcal O_q(t)=\emptyset
  \iff
  \exists\,\bar t:Q\to T,\ \bar t\circ q=t.
\]
This is \(\operatorname{TargetEntropyZero}(q,t)\) iff target descent.

\emph{(3) Target unresolved iff target does not descend.}
The target-unresolved predicate is witnessed by a split pair
\((h,h')\in\mathcal O_q(t)\).  By (2), nonemptiness of
\(\mathcal O_q(t)\) is the negation of target descent.

\emph{(4) Coarse fibers decompose into fine fibers.}
Let \(b\in Q_C\) and \(h\in H\).  If \(q_C(h)=b\), set
\(a:=q_F(h)\).  The factorization \(c\circ q_F=q_C\) gives
\[
  c(a)=c(q_F(h))=q_C(h)=b,
\]
and by construction \(q_F(h)=a\).  Conversely, if there exists
\(a\in Q_F\) with \(c(a)=b\) and \(q_F(h)=a\), then the same
factorization gives
\[
  q_C(h)=c(q_F(h))=c(a)=b.
\]

\emph{(5) Entropy is monotone under coarsening.}
Apply the chain-rule identity to the factorized coarsening pair
\((q_F,q_C)\):
\[
  E(q_C)=E(q_F)\odot m(q_F,q_C).
\]
Combine-monotonicity with \(a:=E(q_F)\) and \(b:=m(q_F,q_C)\) gives
\[
  E(q_F)\le_E E(q_F)\odot m(q_F,q_C)=E(q_C).
\]
Thus \(E(q_F)\le_E E(q_C)\).

\emph{(6) Target entropy as a mixing value.} Under the added hypothesis,
\(m(\langle q,t\rangle,q)=0\) exactly when the two quotients have the same
fibers, that is, when \(q(x)=q(y)\) implies
\((q(x),t(x))=(q(y),t(y))\), the converse being automatic; that is, exactly
when \(t\) is constant on \(q\)-fibers, which by (2) is descent of \(t\) and
\(\operatorname{TargetEntropyZero}(q,t)\). For the fiber-volume functional,
\(m(q',q'')=|H|^{-1}\sum_h\log\bigl(|q''^{-1}(q''h)|/|q'^{-1}(q'h)|\bigr)\) is
a sum of nonnegative terms, since each \(q'\)-fiber lies in a \(q''\)-fiber, and
it vanishes exactly when every term does, that is, when the fibers coincide.
\end{proof}

\subsection*{Layer instantiations}\label{layer-inst:F39}

\begingroup\small
\begin{description}
\item[\emph{Simpson concentration index (ecology).}] \textit{Analogy.}
For a community of \(N\) individuals classified into
species, the access quotient sends each of the \(N\) individuals to its species.  The
Simpson concentration index
\(\lambda=\sum_i p_i^2\) (with \(p_i\) the abundance of
species \(i\)) is the probability that two individuals drawn with
replacement lie in one species fiber; it is not the fiber-volume functional
\(E\) of \Cref{thm:F39} (for the species quotient \(E=\sum_i p_i\log(Np_i)\)),
but, like \(E\), it does not decrease under coarsening.  Taxonomic coarsening (genus,
family) combines species classes; since
\(p_g=\sum_{i\in g}p_i\) and the cross terms of \(p_g^2\) are
nonnegative, \(\sum_g p_g^2\ge\sum_i p_i^2\), so the index is
monotone non-decreasing under coarsening
\citep{Simpson1949Diversity}.

\item[\emph{Burnside orbit-counting in finite group
theory (group theory).}] \textit{Special case.}
For a finite group \(G\) acting on a finite nonempty set \(X\),
the access quotient sends each element to its
\(G\)-orbit. Take any target readout \(t\), for instance \(t(x)=|Gx|\).  The fiber over an orbit \(O\) has
cardinality \(|O|\). For every quotient \(q'\) of the finite nonempty set
\(X\), take \(\mathcal E=[0,\infty)\), \(E(q')=\log(|X|/|\operatorname{im}q'|)\)
and \(\odot=+\); for a coarsening \(q_C=c\circ q_F\), set
\(m(q_F,q_C)=\log(|\operatorname{im}q_F|/|\operatorname{im}q_C|)\). Then
\(m\ge0\), the chain rule and combine-monotonicity of \Fref{F39} hold, and
\(E(q')=0\) exactly when \(q'\) is injective.  This \(E\) is not the
fiber-volume functional of the theorem's last clause; the two agree
when all fibers have the same size. The orbit quotient has entropy
\(\log(|X|/|X/G|)\), with \(|X/G|\) given by
Burnside's lemma as the average number of fixed
points of group elements.  Coarsening to a supergroup
\(G'\supseteq G\) merges orbits into larger fibers,
non-decreasing the average orbit size
\citep{Burnside1897Theory}.

\end{description}
\endgroup


\subsection{Irreducibility / No Short Closed Description [\texorpdfstring{\Fref{F42}}{F42}]}\label{sec:self-reference-limits:irreducibility}

The notion of an object's shortest closed description is the central
object of Kolmogorov complexity and algorithmic information theory;
the standard reference is \citet{LiVitanyi2008Kolmogorov}.

\paragraph{Reading the statement.} The symbols are those of \Cref{thm:F42} below.
\(\mathcal O_t(e,h)\) is the set of pairs in the \(e\)-class of \(h\)
with different target values, and \(\mathcal O_\gamma(e,h)\) the set of
pairs in that class whose compressed states differ after the route
\(r_\gamma\). \(\operatorname{Irreducible}_e(h)\) says that no closed
compression \(e'\) in the admissible class \(\mathcal A\) describes \(h\) more
cheaply than \(e\) does. The predicate does not require \(e\) itself to lie in
\(\mathcal A\) or to be closed for \(h\); it is meant for \(e\in\mathcal A\) with
\(\operatorname{ClosedCompression}(e,h)\), as produced by the existence clause. The class matters: closedness at \(h\) looks only at
the class of \(h\), so if the admissible maps can give \(h\) any value on a
class of its own (for instance, if every map \(H\to E\) is admissible and
\(E\) has at least two elements), irreducibility reduces to minimality of \(\operatorname{cost}(e(h))\) over all of \(E\).

\begin{theorem}[Irreducibility / No Short Closed Description]\label{thm:F42}
\emph{Pointwise form.} Let \(H\) be a set of objects with compression
map \(e:H\to E\), target readout \(t:H\to T\), and route updates
\(r_\gamma:H\to H\) for \(\gamma\in\Gamma\). Let \(\mathcal A\) be a declared
class of admissible compressions \(H\to E\), and let
\(\operatorname{cost}:E\to C\) take values in a set \(C\) with a relation \(<\) (no order axioms are used; the existence clause uses only well-foundedness). For \(h\in H\) put
\(\mathcal O_t(e,h):=\{(x,x')\in H^2:e(x)=e(x')=e(h),\ t(x)\ne t(x')\}\) and
\(\mathcal O_\gamma(e,h):=\{(x,x')\in H^2:e(x)=e(x')=e(h),\ e(r_\gamma x)\ne
e(r_\gamma x')\}\). A compression \(e\) is closed for \(h\), written
\(\operatorname{ClosedCompression}(e,h)\), when \(t\) and every
\(e\circ r_\gamma\) are constant on the class \(e^{-1}(e(h))\). Define
\[
  \operatorname{Irreducible}_e(h)
  :\Longleftrightarrow
  \neg\exists e'\in\mathcal A,\
  \operatorname{ClosedCompression}(e',h)
  \ \wedge\
  \operatorname{cost}(e'(h))<\operatorname{cost}(e(h)).
\]
Then
\[
  \operatorname{ClosedCompression}(e,h)
  \iff
  \mathcal O_t(e,h)=\emptyset
  \ \wedge\
  \forall\gamma\in\Gamma,\ \mathcal O_\gamma(e,h)=\emptyset.
\]
By the displayed equivalence, \(\operatorname{Irreducible}_e(h)\) holds
exactly when every \(e'\in\mathcal A\) with
\(\operatorname{cost}(e'(h))<\operatorname{cost}(e(h))\) has
\(\mathcal O_t(e',h)\ne\emptyset\) or \(\mathcal O_\gamma(e',h)\ne\emptyset\)
for some \(\gamma\). Moreover, if \(<\) is well founded on \(C\) and some compression
in \(\mathcal A\) is closed for \(h\), then there is \(e^\ast\in\mathcal A\)
that is closed for \(h\) and satisfies
\(\operatorname{Irreducible}_{e^\ast}(h)\). If for every \(v\in E\) the class \(\mathcal A\) contains a map \(e'\) with
\(e'(h)=v\) and \(e'(x)\ne v\) for every \(x\ne h\) (for instance, if
\(\mathcal A\) contains every map \(H\to E\) and \(E\) has at least two
elements), then
\(\operatorname{Irreducible}_e(h)\) holds exactly when no \(v\in E\) has
\(\operatorname{cost}(v)<\operatorname{cost}(e(h))\). For such a class
\(\mathcal A\) the pointwise notion depends on neither \(t\) nor the routes
\(r_\gamma\); it constrains the target only when \(\mathcal A\) excludes maps
isolating \(h\). The global notion below requires closure on every class, so an
isolating map is globally closed only when \(t\) and every \(r_\gamma\) respect
all its classes.
\emph{Global form.} Let \(b:H\to B\) be a map. Separately declare an admissible
class \(\mathcal A_{\mathrm{glob}}\) of descriptions \(e:H\to E'\), allowing the
codomain \(E'\) to vary, and a cost \(c\) on these descriptions and on \(b\).
Call \(e\) \emph{globally closed over} \(b\) when \(e\) descends through \(b\),
\(t\) descends through \(e\), and \(e\circ r_\gamma\) descends through \(e\) for
every \(\gamma\). For surjective \(b\) and \(e\), this holds exactly when the sets
\(\{(x,x'):b(x)=b(x'),\ e(x)\ne e(x')\}\) and
\(\{(x,x'):e(x)=e(x'),\ t(x)\ne t(x')\}\) and, for every \(\gamma\),
\(\{(x,x'):e(x)=e(x'),\ e(r_\gamma x)\ne e(r_\gamma x')\}\) are empty. Define
\(b\) to be \emph{globally irreducible} when no \(e\in\mathcal A_{\mathrm{glob}}\)
globally closed over \(b\) satisfies \(c(e)<c(b)\). If \(c\) takes values in a
well-founded strict order and some description in \(\mathcal A_{\mathrm{glob}}\)
is globally closed over \(b\), then some \(e^\ast\in\mathcal A_{\mathrm{glob}}\)
globally closed over \(b\) has no description in \(\mathcal A_{\mathrm{glob}}\)
globally closed over \(b\) of smaller cost, and \(e^\ast\) is globally
irreducible (with \(e^\ast\) in place of \(b\)).
\end{theorem}
\begin{proof}
Assume the named maps and the declared cost relation. A compression is
closed exactly when no target or route update distinguishes two
representatives of the same compressed description. In symbols,
\[
  \mathcal O_t(e,h)=
  \{(x,x'):e(x)=e(x')=e(h),\ t(x)\ne t(x')\},
\]
and similarly \(\mathcal O_\gamma(e,h)\) is formed using
\(e(r_\gamma x)\). Hence closedness is precisely the simultaneous
emptiness of these obstruction sets. Substituting this for each cheaper \(e'\in\mathcal A\) in the definition of
\(\operatorname{Irreducible}_e(h)\) gives the reformulation after the display:
a shorter closed \(e'\) carries all declared readouts with smaller cost, so
\(h\) is reducible exactly when some cheaper admissible candidate has all its
obstruction sets empty.
If \(\operatorname{Irreducible}_e(h)\) holds, no admissible closed compression
has cost below \(\operatorname{cost}(e(h))\). If moreover \(e\in\mathcal A\)
is closed at \(h\), then \(e\) is minimal among admissible closed compressions.
For the existence clause, the costs \(\operatorname{cost}(e'(h))\) of the
closed compressions \(e'\in\mathcal A\) at \(h\) form a nonempty subset of
\(C\). Since \(<\) is well founded, this subset has a \(<\)-minimal element
\(\operatorname{cost}(e^\ast(h))\) with \(e^\ast\in\mathcal A\) closed at
\(h\), and no closed compression in \(\mathcal A\) at \(h\) has a cost below
it, which is \(\operatorname{Irreducible}_{e^\ast}(h)\).
For the last clause, let \(v\in E\) and take \(e'\in\mathcal A\) with
\(e'(h)=v\) and \(e'(x)\ne v\) for \(x\ne h\); when \(\mathcal A\) contains
every map and \(E\) has two elements, choose \(w\ne v\) and set
\(e'(h)=v\) and \(e'(x)=w\) for \(x\ne h\). The map \(e'\) has
\(\{h\}\) as the \(e'\)-class of \(h\), so both obstruction sets are empty and
\(e'\) is closed for \(h\). Hence a \(v\) with
\(\operatorname{cost}(v)<\operatorname{cost}(e(h))\) gives a cheaper closed
\(e'\in\mathcal A\); conversely, a cheaper closed \(e'\) gives such a \(v\),
namely \(e'(h)\). For the global clause, each descent condition is the
emptiness of the corresponding set by the descent criterion, using
surjectivity of \(b\) and \(e\). An isolating map is globally closed only if
\(t\) and every \(r_\gamma\) respect its classes on all of \(H\). For the global
minimum, take a closed description of minimal cost; any description globally
closed over \(e^\ast\) descends through \(e^\ast\), hence through \(b\), so it is
closed over \(b\) and cannot cost less.
\end{proof}

\subsection*{Layer instantiations}\label{layer-inst:F42}

\begingroup\small
\begin{description}
\item[\emph{Abel-Ruffini-Galois unsolvability of the
quintic (algebra).}] \textit{Analogy.}
For a quintic over \(\mathbb Q\) with Galois group \(S_5\), no
root is a nested-radical expression in its coefficients, by Galois's solvability criterion; Abel proved the corresponding statement for the general quintic \citep{Abel1824Quintic,Galois1846Memoire}. To instantiate \Fref{F42}, one must
specify \(H,E,t,\mathcal A\) and establish that every compression
in \(\mathcal A\) closed at the chosen polynomial would yield such
a radical expression for the target root.

\end{description}
\endgroup


\subsection{Continuum Emergence [\texorpdfstring{\Fref{F43}}{F43}]}\label{sec:self-reference-limits:continuum-emergence}

\paragraph{Reading the statement.} \Cref{thm:F43} below does not
construct a limit: the candidate limit \(L\), the projections \(p_i\) and
the limit quotient \(\ell\) are given, and no topology or convergence
beyond the declared predicate \(\operatorname{Vanishes}\) on the
residuals is used.  \(\operatorname{Cofinal}(I)\) says that every
declared finite test is covered at some index; it is a covering
condition, not order-theoretic cofinality.  Coherence, cofinality and
vanishing enter \(\operatorname{StableLimitQuotient}(L)\) only as
assumed conjuncts and are carried unchanged into the conclusion.  The
content of the theorem is the existence and uniqueness of a readout
\(\bar\ell\) on the limit quotient that agrees with every stage readout:
the equations \(\phi_{ij}\circ p_j=p_i\) and \(r_i\circ\phi_{ij}=r_j\),
together with directedness, make the stage readouts agree, and
\(\bar\ell\) exists when their common value descends through \(\ell\).

\begin{theorem}[Continuum Emergence]\label{thm:F43}
Let \((I,\le)\) be a nonempty directed index set, let \(A_i\) be a discrete
approximation at index \(i\), and let
\(\phi_{ij}:A_j\to A_i\) be comparison maps for \(i\le j\). Let
\(p_i:L\to A_i\) be projections from a candidate limit \(L\), let
\(\xi_i\in R_i\) be residuals, let
\(\ell:L\to Q_\infty\) be a surjective limit quotient, let
\(r_i:A_i\to L_{\mathrm{rec}}\) be stage readouts, and let
\(t:L\to L_{\mathrm{rec}}\) be a target readout. Let \(\mathcal T\) be a set of
declared finite tests with a covering relation
\(\operatorname{cov}\subseteq I\times\mathcal T\), write
\(\operatorname{Cofinal}(I)\) (a covering condition, not order-theoretic
cofinality) \(:\Longleftrightarrow\forall\tau\in\mathcal T\ \exists i\in I\ \operatorname{cov}(i,\tau)\),
and let \(\operatorname{Vanishes}\) be a declared predicate on residual
families \((\xi_i)_{i\in I}\). Write \(\operatorname{TowerCoherent}(\phi_{ij})\)
when \(\phi_{ii}=\mathrm{id}\) and \(\phi_{ik}=\phi_{ij}\circ\phi_{jk}\) for
\(i\le j\le k\), and call \(\bar\ell:Q_\infty\to L_{\mathrm{rec}}\) compatible with
all \(p_i\) when \(\bar\ell(\ell(x))=r_i(p_i(x))\) for all \(i\in I\) and
\(x\in L\). Write
\begin{align*}
  \operatorname{StableLimitQuotient}(L)
  :\iff\
  &\operatorname{TowerCoherent}(\phi_{ij})
  \ \wedge\
  \operatorname{Cofinal}(I)\\
  &\wedge\
  \operatorname{Vanishes}(\xi_i)
  \ \wedge\
  \exists\,\bar\ell:Q_\infty\to L_{\mathrm{rec}}\ \text{ compatible with all }p_i .
\end{align*}
Suppose that the tower is coherent and cofinal, that the residuals vanish,
that \(\phi_{ij}\circ p_j=p_i\) and \(r_i\circ\phi_{ij}=r_j\) for \(i\le j\),
and that for some index \(b\) the target readout satisfies
\(t=r_b\circ p_b\) and is constant on the fibers of \(\ell\). Then there is exactly one \(\bar\ell:Q_\infty\to L_{\mathrm{rec}}\) compatible
with all \(p_i\), namely \(\bar\ell(\ell(x))=t(x)\); with the assumed coherence,
cofinality and vanishing, \(\operatorname{StableLimitQuotient}(L)\) holds. More generally, assume only that \(\phi_{ij}\circ p_j=p_i\) and
\(r_i\circ\phi_{ij}=r_j\) for \(i\le j\). Then \(r_i\circ p_i=r_j\circ p_j\) for all
\(i,j\in I\); for every index \(b\), a map \(\bar\ell\) compatible with all \(p_i\)
exists if and only if \(r_b\circ p_b\) is constant on the fibers of \(\ell\), and
it is then unique; hence \(\operatorname{StableLimitQuotient}(L)\) holds if and
only if, in addition, the tower is coherent and cofinal and the residuals
vanish.
\end{theorem}
\begin{proof}
Fix \(x\in L\) and \(i\in I\). Since \(I\) is directed, choose \(k\) with
\(b\le k\) and \(i\le k\). Using \(p_b=\phi_{bk}\circ p_k\) and
\(r_b\circ\phi_{bk}=r_k\),
\[
  t(x)=r_b(p_b(x))=r_b(\phi_{bk}(p_k(x)))=r_k(p_k(x)),
\]
and in the same way
\(r_i(p_i(x))=r_i(\phi_{ik}(p_k(x)))=r_k(p_k(x))\). Hence
\(r_i(p_i(x))=t(x)\) for every \(i\). Since \(t\) is constant on the fibers
of \(\ell\) and \(\ell\) is surjective, \(\bar\ell(\ell(x)):=t(x)\) defines a
map \(\bar\ell:Q_\infty\to L_{\mathrm{rec}}\), and
\(\bar\ell(\ell(x))=r_i(p_i(x))\) for all \(i\) and \(x\), so \(\bar\ell\) is
compatible with all \(p_i\). Together with the coherence, cofinality and
vanishing hypotheses, the four conditions defining
\(\operatorname{StableLimitQuotient}(L)\) hold.
For the general clause, the same computation with a common upper bound \(k\)
of \(i\) and \(j\) gives \(r_i\circ p_i=r_k\circ p_k=r_j\circ p_j\). If
\(\bar\ell\) is compatible with all \(p_i\), then
\(r_b(p_b(x))=\bar\ell(\ell(x))\) is constant on the fibers of \(\ell\);
conversely, if \(r_b\circ p_b\) is constant on those fibers, the argument above
with \(t:=r_b\circ p_b\) gives a compatible \(\bar\ell\). Since \(\ell\) is
surjective, \(\bar\ell\) is determined by \(\bar\ell(\ell(x))=r_b(p_b(x))\).
\end{proof}

The result supplies a readout on the specified quotient of a candidate limit
that agrees with every stage readout.

\subsection*{Layer instantiations}\label{layer-inst:F43}

\begingroup\small
\begin{description}
\item[\emph{Dedekind cut construction of the reals
(analysis).}] \textit{Analogy.}
The discrete approximations are rational intervals
indexed by precision; comparison maps are
inclusions of finer intervals into coarser ones.
A Dedekind cut is a non-empty proper downward-closed
subset \(A\subsetneq\mathbb Q\) without greatest element.  The candidate limit is the set of cuts;
projections send a cut to its rational interval
approximant; the surjective limit quotient sends
a cut to its real value, and tower coherence is
nested-approximation consistency.  The analogy
is informal: the cut construction is not
presented here as a tower satisfying the
hypotheses of \Fref{F43} \citep{Dedekind1872Stetigkeit}.

\end{description}
\endgroup


\subsection{Criticality / Phase Boundary [\texorpdfstring{\Fref{F44}}{F44}]}\label{sec:self-reference-limits:criticality}

The motivating substrate analogy is the theory of phase transitions and
critical phenomena; the standard textbook treatment is
\citet{Stanley1971PhaseTransitions}.  The layer-agnostic statement
below records the descent-and-margin shape that the substrate
realizes.

\paragraph{Reading the statement.} In \Cref{thm:F44} below, the three
displayed lines define phase transition, phase boundary and criticality from
the declared data, and clause (i) is the ordinary quotient-descent criterion. Neighborhoods are indexed by \(W\in\mathcal W\), and \(U\) is the non-uniform-response
predicate. For a general relation \(N\) the predicate is only as informative
as \(N\): if \(\mathcal W\) is empty, every formed control is a phase boundary,
and an index \(W\) with no \(y\) satisfying \(N(x,W,y)\) makes \(x\) locally
phase-constant. The relation of (iv) excludes both cases. Clause (iii) asks that, at each step and for every \(W\), some
neighbour of \(x_k\) have the phase of \(x_{k+1}\); with the neighbourhood
relation of (iv) this says that \(x_k\) lies in the closure of the phase class
of \(x_{k+1}\), and (iv) is the form that applies to continuous paths.
The relation of (iv) is chosen so that the empty open set does not
witness local constancy, as it would for the relation
\(x\in W\wedge y\in W\).

\begin{theorem}[Criticality / Phase Boundary]\label{thm:F44}
Let \(\mathcal C\) be a control space with surjective quotient
\(c:\mathcal C\to C\), phase readout \(\varphi:\mathcal C\to\Phi\),
formedness predicate \(F:\mathcal C\to\mathrm{Prop}\), and
neighborhood relation \(N\subseteq\mathcal C\times\mathcal W\times\mathcal C\) over a
declared set \(\mathcal W\) of neighbourhood indices (in (iv), \(\mathcal W\) is
the set of open subsets of \(\mathcal C\)). Let
\(M,U,R:\mathcal C\to\mathrm{Prop}\) denote, respectively,
margin-vanishing, non-uniform response, and
residuals-statused. For a readout \(r:\mathcal C\to V\) let
\(\mathcal O_c(r)=\{(x,y):c(x)=c(y)\ \text{and}\ r(x)\ne r(y)\}\), and for
\(x,x_-,x_+\in\mathcal C\) define
\begin{align*}
  \operatorname{PhaseTransition}(x_-,x_+)
  &:\Longleftrightarrow \varphi(x_-)\ne\varphi(x_+),\\
  \operatorname{PhaseBoundary}(x)
  &:\Longleftrightarrow F(x)\ \wedge\
    \neg\exists W,\ \forall y,\ N(x,W,y)\Rightarrow \varphi(y)=\varphi(x),\\
  \operatorname{Critical}(x)
  &:\Longleftrightarrow \operatorname{PhaseBoundary}(x)\wedge (M(x)\vee U(x))\wedge R(x).
\end{align*}
Then:
\begin{enumerate}[label=(\roman*)]
\item \(r\) descends through \(c\), that is \(r=\bar r\circ c\) for some
  \(\bar r:C\to V\), if and only if \(\mathcal O_c(r)=\emptyset\);
\item if \(\varphi\), \(F\), \(M\), \(U\) and \(R\) are constant on the fibers
  of \(c\), and \(N\) respects \(c\) (\(N(x,W,y)\iff N(x',W,y)\) whenever
  \(c(x)=c(x')\)), then \(\operatorname{PhaseBoundary}\) and
  \(\operatorname{Critical}\) are constant on the fibers of \(c\), so by (i)
  they descend through \(c\): whether a control is critical depends only on
  its class;
\item if \(x_0,\dots,x_n\) is a finite list of controls with \(F(x_k)\) for
  every \(k\) and, for every \(k<n\) and every \(W\), some \(y\) with
  \(N(x_k,W,y)\) has \(\varphi(y)=\varphi(x_{k+1})\) (for instance,
  \(N(x_k,W,x_{k+1})\) for every \(W\), as when \(N\) is a neighbourhood graph
  whose relation does not depend on \(W\); with the topological relation \(N(x,W,y):\Leftrightarrow(x\in W\Rightarrow y\in W)\), \(W\) open, used in (iv), the condition says exactly that \(x_k\) lies in the closure of the phase class of \(x_{k+1}\)), and
  \(\varphi(x_0)\ne\varphi(x_n)\), then some \(k<n\) has
  \(\varphi(x_k)\ne\varphi(x_{k+1})\) and
  \(\operatorname{PhaseBoundary}(x_k)\): a phase transition along a path of
  neighbouring formed controls crosses a phase boundary;
\item if \(\mathcal C\) is a topological space, \(W\) ranges over its open
  sets and \(N(x,W,y)\) means that \(y\in W\) whenever \(x\in W\), and
  \(\gamma:[0,1]\to\mathcal C\) is continuous with \(F(\gamma(s))\) for every
  \(s\) and \(\varphi(\gamma(0))\ne\varphi(\gamma(1))\), then
  \(\operatorname{PhaseBoundary}(\gamma(s))\) for some \(s\in[0,1]\): a phase
  transition along a continuous path of formed controls crosses a phase
  boundary. With this \(N\), \(\operatorname{PhaseBoundary}(x)\) holds exactly
  when \(F(x)\) holds, \(\varphi\) is not constant on \(\mathcal C\), and
  \(\varphi\) is constant on no open neighbourhood of \(x\).
\end{enumerate}
\end{theorem}
\begin{proof}
Since \(c\) is surjective, (i) is the descent criterion of
\Cref{sec:apparatus}.  The defined predicates sort a formed control \(x\) into two cases:
\[
  \begin{cases}
  \exists W\,\forall y,\ N(x,W,y)\Rightarrow\varphi(y)=\varphi(x),
    &\text{\(x\) is locally phase-constant,}\\
  \neg\exists W\,\forall y,\ N(x,W,y)\Rightarrow\varphi(y)=\varphi(x),
    &\text{\(x\) lies on a phase boundary,}
  \end{cases}
\]
and a critical control is a boundary control at which the separating margin
vanishes or the response is non-uniform, with its residuals statused.
For (ii), let \(c(x)=c(x')\). Then \(F(x)=F(x')\) and
\(\varphi(x)=\varphi(x')\), and for every \(W\) and \(y\),
\(N(x,W,y)\iff N(x',W,y)\); so \(x\) is locally phase-constant exactly when
\(x'\) is, and \(\operatorname{PhaseBoundary}(x)\iff
\operatorname{PhaseBoundary}(x')\). With \(M(x)=M(x')\), \(U(x)=U(x')\) and
\(R(x)=R(x')\), also \(\operatorname{Critical}(x)\iff
\operatorname{Critical}(x')\).
(iii) Since \(\varphi(x_0)\ne\varphi(x_n)\), some \(k<n\) has
\(\varphi(x_{k+1})\ne\varphi(x_k)\). For every \(W\), the hypothesis gives
\(y\) with \(N(x_k,W,y)\) and \(\varphi(y)=\varphi(x_{k+1})\ne\varphi(x_k)\), so
no \(W\) makes \(\varphi\) constant at \(x_k\); with \(F(x_k)\) this is
\(\operatorname{PhaseBoundary}(x_k)\).
(iv) Suppose no \(\gamma(s)\) is a boundary point. Since every
\(\gamma(s)\) is formed, for each \(s\) there is an open \(W_s\) with
\(\varphi(y)=\varphi(\gamma(s))\) for every \(y\) such that
\(\gamma(s)\in W_s\Rightarrow y\in W_s\). If \(\gamma(s)\notin W_s\), every
\(y\) qualifies, so \(\varphi\) is constant on \(\mathcal C\), contradicting
\(\varphi(\gamma(0))\ne\varphi(\gamma(1))\). Hence \(\gamma(s)\in W_s\) and
\(\varphi\) is constant on \(W_s\), so \(\varphi\circ\gamma\) is constant on
the open neighbourhood \(\gamma^{-1}(W_s)\) of \(s\). A locally constant map
on the connected interval \([0,1]\) is constant, so
\(\varphi(\gamma(0))=\varphi(\gamma(1))\), a contradiction.
For the last sentence of (iv): an open \(W\not\ni x\) relates \(x\) to every
\(y\), so it witnesses local constancy exactly when \(\varphi\) is constant on
\(\mathcal C\); an open \(W\ni x\) does so exactly when \(\varphi\) is constant on
\(W\). Since \(\mathcal C\) is an open neighbourhood of \(x\),
\(\operatorname{PhaseBoundary}(x)\) holds exactly when \(F(x)\) holds and
\(\varphi\) is constant on no open neighbourhood of \(x\), which implies that
\(\varphi\) is not constant on \(\mathcal C\).
\end{proof}

\subsection*{Layer instantiations}\label{layer-inst:F44}

\begingroup\small
\begin{description}
\item[\emph{Allee effect threshold (population
ecology).}] \textit{Analogy.}
For a species population \(\dot N=N\,r(N)\), the
control space is population size; the phase readout
records ``extinction'' vs ``persistence''.  The Allee
effect makes the per-capita growth \(r(N)\) negative
below an Allee threshold \(N_A\), producing a phase
boundary at \(N=N_A\): the per-capita growth rate
\(r(N_A)\) vanishes, response to environmental perturbations is
non-uniform, and the residuals-statused predicate \(R\) would have
to be supplied as a declared status of demographic
noise, which the analogy does not provide
\citep{Allee1931Aggregations}.

\item[\emph{Thom cusp catastrophe (catastrophe
theory).}] \textit{Analogy.}
The cusp catastrophe is the family of potentials
\(V_a(x)=x^4/4+a_1 x^2/2+a_2 x\) parameterised by
\((a_1,a_2)\in\mathbb R^2\); the phase readout
records the number of stable equilibria.  The
bifurcation set \(4a_1^3+27a_2^2=0\) is the phase boundary: on its
fold branches one stable equilibrium merges with the unstable one,
all three equilibria merge at the cusp point \(a=0\), the energy
margin vanishes, the response is
non-uniform, and the residual smooth-perturbation
status is classified by the elementary catastrophe
normal form \citep{Thom1972Stabilite}.

\item[\emph{{NBER} business-cycle peak/trough
(macroeconomics).}] \textit{Analogy.}
The control space is macroeconomic time series of
GDP, employment, industrial production, and personal
income; the phase readout records expansion vs
recession.  NBER-defined peaks and troughs are
phase boundaries: leading-indicator marginal growth
vanishes, sectoral responses are non-uniform, and
the residual status is handled by the NBER Business
Cycle Dating Committee's revision protocol over
noisy indicators \citep{BurnsMitchell1946Measuring}.

\item[\emph{Hopf bifurcation onset of oscillation
(nonlinear dynamics).}] \textit{Analogy.}
For a parameter-dependent ODE
\(\dot{\mathbf x}=f(\mathbf x;\mu)\), the control
space is \(\mu\); the phase readout records ``stable
fixed point'' vs ``limit cycle''.  A Hopf bifurcation
at \(\mu=\mu_H\) occurs when a complex-conjugate
eigenvalue pair of the Jacobian crosses the
imaginary axis: the real part vanishes, in the supercritical case
oscillation amplitude grows as \(\sqrt{\mu-\mu_H}\) (a subcritical
bifurcation instead produces a jump to a distant attractor), and the
residual perturbation is statused by smooth
normal-form coefficients \citep{Hopf1942Abzweigung}.
\end{description}
\endgroup


\subsection{Locality / Domain of Dependence [\texorpdfstring{\Fref{F45}}{F45}]}\label{sec:self-reference-limits:locality}

The classical relativistic analogy for a local quotient with a
domain-of-dependence structure is the causal structure of
Lorentzian spacetime; the standard reference for the
domain-of-dependence formalism is \citet[Chapter~8]{Wald1984GR}.

\begin{theorem}[Locality / Domain of Dependence]\label{thm:F45}
Let \(H\) be a history set with local quotient \(\ell:H\to L\),
target readout \(t:H\to T\), signature \(\sigma:H\to\Sigma\), update
map \(u:H\to H'\), and target quotient \(q':H'\to Q'\). Assume
\(\ell\) is surjective. Let \(<\) be a declared strict-smaller relation on a
declared class of local quotients of \(H\) that contains \(\ell\); in
\(\operatorname{MinimalDomain}\), \(\ell'\) ranges over that class. Write
\(\operatorname{LocalToDomain}(\ell,t)\) when \(t\) factors through \(\ell\),
that is \(t=\bar t\circ\ell\) for some \(\bar t:L\to T\);
\(\operatorname{TargetFamilyLocal}(\ell,\sigma)\) when \(\sigma\) factors
through \(\ell\); and \(\operatorname{DynamicLocality}(\ell,u,q')\) when
\(q'\circ u\) factors through \(\ell\). Define
\[
  \mathcal O_\ell(t)=
  \{(h,h'):\ell(h)=\ell(h')\ \text{and}\ t(h)\ne t(h')\},
\]
and define
\[
  \operatorname{MinimalDomain}(\ell,\sigma)
  :\Longleftrightarrow \operatorname{TargetFamilyLocal}(\ell,\sigma)\ \wedge\
  \neg\exists \ell',\ \ell'\!<\ell\ \wedge\
  \operatorname{TargetFamilyLocal}(\ell',\sigma).
\]
Then
\begin{align*}
  \operatorname{LocalToDomain}(\ell,t)
  &\iff \mathcal O_\ell(t)=\emptyset,\\
  \operatorname{TargetFamilyLocal}(\ell,\sigma)
  &\iff
  \forall a,b\in H,\ \ell(a)=\ell(b)\Rightarrow \sigma(a)=\sigma(b),\\
  \operatorname{DynamicLocality}(\ell,u,q')
  &\iff
  \forall a,b\in H,\ \ell(a)=\ell(b)\Rightarrow q'(u(a))=q'(u(b)).
\end{align*}
Moreover, the local repair \(\langle\ell,t\rangle:H\to L\times T\) satisfies
\(\operatorname{LocalToDomain}(\langle\ell,t\rangle,t)\). If \(\ell'<\ell\)
means that \(\ell'\) factors through \(\ell\) and \(\ell\) does not factor
through \(\ell'\), the declared class contains \(\sigma\), and \(H\) is
nonempty, then \(\operatorname{MinimalDomain}(\ell,\sigma)\) holds if and only if \(\ell\) and
\(\sigma\) have the same fibers.
\end{theorem}
\begin{proof}
Since \(\ell\) is surjective, the three displayed equivalences are the
descent criterion of \Cref{sec:apparatus} for \(\ell\) and the readouts
\(t\), \(\sigma\) and \(q'\circ u\).  For the local repair, \(t=\mathrm{pr}_2\circ\langle\ell,t\rangle\), so \(t\)
factors through \(\langle\ell,t\rangle\). For the last clause, note that for
nonempty \(H\) two maps on \(H\) with the same fibers factor through each other
(for empty \(H\) a map into an empty codomain need not exist). Suppose
\(\operatorname{MinimalDomain}(\ell,\sigma)\). Then \(\sigma\) factors through
\(\ell\); if \(\ell\) did not factor through \(\sigma\), then \(\sigma<\ell\)
and \(\sigma\) factors through itself, contradicting minimality, so \(\ell\)
and \(\sigma\) have the same fibers. Conversely, if they have the same fibers
and \(\sigma\) factors through some \(\ell'\) with \(\ell'<\ell\), then
\(\ell\), which factors through \(\sigma\), factors through \(\ell'\),
contradicting \(\ell'<\ell\).
\end{proof}

\subsection*{Layer instantiations}\label{layer-inst:F45}

\begingroup\small
\begin{description}
\item[\emph{Markov property in stochastic processes
(probability theory).}] \textit{Special case.}
For a time-homogeneous Markov chain \((X_t)\) on a countable
state space \(S\) with transition matrix \(P\), let \(H\) be the set of
state sequences \(\omega=(x_0,\dots,x_k)\) of positive probability,
\(\ell(\omega)=x_k\) and \(L=\ell(H)\), so \(\ell\) is surjective. For an
integer \(\Delta\ge1\), let \(t(\omega)\) be the law of \(X_{k+\Delta}\) given
\((X_0,\dots,X_k)=\omega\). The Markov property says
\(t(\omega)=P^{\Delta}(x_k,\cdot)\), which implies
\(\operatorname{LocalToDomain}(\ell,t)\) with \(\bar t(x)=P^{\Delta}(x,\cdot)\). Take \(\sigma=t\), all quotients of \(H\)
as the declared class, and strict coarsening as \(<\). Then \(\ell\)
is sufficient and is minimal exactly when distinct states in \(L\)
have distinct \(\Delta\)-step rows; otherwise equality of those rows
gives a coarser sufficient quotient. This descent condition concerns the \(\Delta\)-step law of the unlumped state; lumpability asks that the lumped process itself be Markov, and the two conditions can coincide \citep{Markov1906Investigation}.

\item[\emph{Sheaf locality on a topological space
(algebraic geometry).}] \textit{Special case.}
For a sheaf \(\mathcal F\) on a topological space
\(X\), the history set consists of pairs
\((U,s)\) of an open subset and a section; the
local quotient sends \((U,s)\) to \(U\) together with its family of
germs \((s_x)_{x\in U}\), \(s_x\in\mathcal F_x\), corestricted to its image so that \(\ell\) is surjective.  The target
readout is the section.  The sheaf condition
states that two sections agreeing on all germs
coincide, so the section descends through the germ-family quotient (gluing is not needed, since every history is an actual section).  Take \(\sigma=t\), all quotients of \(H\) as
the declared class and strict coarsening as \(<\). Since the family of
germs and the section determine each other, \(\ell\) and the readout have
the same fibers, so \(\ell\) is a minimal domain by the last clause of
\Cref{thm:F45} \citep{Grothendieck1957Tohoku}.

\item[\emph{Cellular automaton local-update rule
(theoretical computer science).}] \textit{Special case.}
For a cellular automaton on \(\mathbb Z^d\) with
finite state space, the history set is global
configurations; the local quotient sends a
configuration to its restriction to a radius-\(r\)
neighbourhood of a cell.  The local update rule
computes the next state from the radius-\(r\)
neighbourhood; the next-state readout descends
through the neighbourhood quotient, so the
radius-\(r\) neighbourhood is a sufficient domain.
Take \(\sigma\) the next-state readout, all quotients of the history set as the declared class and strict coarsening as \(<\); by the last clause of \Cref{thm:F45}, a local quotient is minimal exactly when it has the fibers of \(\sigma\). The minimal quotient identifies neighbourhood
patterns with the same next state; for an XOR rule it
is the parity of the neighbourhood
\citep{vonNeumann1966SelfReproducing}.

\item[\emph{Lamport happens-before relation
(distributed systems).}] \textit{Special case.}
For a distributed system of asynchronous processes exchanging
messages, fix an event set containing at least three events. The history set consists of ordered pairs \((e_1,e_2)\) of distinct events; the local quotient \(\ell\) sends such a pair to its pair of causal pasts \((\downarrow e_1,\downarrow e_2)\), corestricted to its image, with \(\downarrow e:=\{f:f\to e\}\cup\{e\}\), where \(\to\) is Lamport's relation, the smallest transitive relation in which events within one process are ordered by occurrence and every message-send precedes its matching receive. Let \(\sigma(e_1,e_2)\) be the verdict \(e_1\to e_2\), \(e_2\to e_1\) or \(e_1\parallel e_2\) (concurrent). Since \(e_1\to e_2\) exactly when \(e_1\in\downarrow e_2\), the verdict descends through \(\ell\). Take all quotients of the history set as the declared class and strict coarsening as \(<\). Because \(e\) is the greatest element of \(\downarrow e\), \(\ell\) is injective, while \(\sigma\) takes at most three values on at least six pairs; so by the last clause of \Cref{thm:F45} \(\ell\) is sufficient but not minimal, and the minimal domain is the verdict itself \citep{Lamport1978HappensBefore}.

\item[\emph{Voronoi cell as nearest-neighbour
domain (computational geometry).}] \textit{Special case.}
For a nonempty finite indexed set of distinct seed points
\(P=\{p_1,\ldots,p_m\}\subset\mathbb R^d\), the history set is the ambient
space; the local quotient sends each point to the
index of its nearest seed, ties broken by the smallest index.  The local
domain is the cell \(\{x:\text{the tie-broken nearest seed is }p_i\}\),
contained in the Voronoi cell
\(V_i=\{x:\|x-p_i\|\le\|x-p_j\|\,\forall j\}\), and
the nearest-neighbour classification descends
through this quotient.  On interior points of these cells the verdict is
robust to small perturbations. At distance ties the declared rule gives a
unique verdict, although it need not be robust to perturbations.
Take \(\sigma\) the verdict readout, all quotients of the history set as the declared class and strict coarsening as \(<\); by the last clause of \Cref{thm:F45}, a local quotient is minimal exactly when it has the fibers of \(\sigma\). The minimal domain is the partition into the tie-broken cells
\citep{Voronoi1908Nouvelles}.
\end{description}
\endgroup

\subsection{Law--State Split [\texorpdfstring{\Fref{F46}}{F46}]}\label{sec:self-reference-limits:law-state-split}

\paragraph{Reading the statement.} The symbols are those of \Cref{thm:F46} below.
By clause (ii), for a readout that descends the split does not depend on
\(\mu\): a law readout is a readout constant on \(\mathcal X\), and \(\mu\) enters
only through descent, uniqueness of \(\bar r\), and clauses (iii) and (iv).
Taking every closure lawful and \(c\sim d:\Longleftrightarrow\mu(c)=\mu(d)\)
in \Cref{thm:F26}, a readout that descends through \(\mu\) is a moduli
readout, \(\operatorname{LawReadout}\) is \(\operatorname{NecessaryReadout}\),
and for a descending readout \(\operatorname{StateDependent}\) is
\(\operatorname{ContingentReadout}\). Clause~(i) below is the first half of
clause~(i) of \Cref{thm:F26} with uniqueness of \(\bar r\) added; clause~(ii)
contains clause~(ii) there; clause~(iv) is clause~(iv) there with \(\sim'\)
the kernel of \(\mu_2\); and clause~(iii) evaluates the descended readout at
the selected modulus \(m_\ast\).

\begin{theorem}[Law--State Split]\label{thm:F46}
Let \(\mathcal X\) be a closure space with surjective moduli quotient
\(\mu:\mathcal X\to M\) and readout \(r:\mathcal X\to V\). Let
\(m_\ast\in M\) be a selected modulus, and let
\(\tau:\mathcal X\to\mathcal X_2\) be a transport into a second closure
space with moduli quotient \(\mu_2:\mathcal X_2\to M_2\). The readout
\(r\) \emph{descends} through \(\mu\) when \(r=\bar r\circ\mu\) for some
\(\bar r:M\to V\), its \emph{descended readout}. For a readout that descends,
define
\begin{align*}
  \operatorname{LawReadout}(r)
  &:\Longleftrightarrow
  \forall m,m'\in M,\ \bar r(m)=\bar r(m'),\\
  \operatorname{StateDependent}(r)
  &:\Longleftrightarrow \neg\operatorname{LawReadout}(r).
\end{align*}
Then:
\begin{enumerate}[label=(\roman*)]
\item \(r\) descends through \(\mu\) if and only if no pair \(x,y\) has
  \(\mu(x)=\mu(y)\) and \(r(x)\ne r(y)\), and then its descended readout
  \(\bar r\) is unique;
\item if \(r\) descends, \(\operatorname{LawReadout}(r)\) holds if and only if no
  pair \(m,m'\in M\) has \(\bar r(m)\ne\bar r(m')\), if and only if \(r\) is
  constant on \(\mathcal X\), and exactly one of
  \(\operatorname{LawReadout}(r)\) and \(\operatorname{StateDependent}(r)\)
  holds; moreover \(\operatorname{StateDependent}(r)\) holds if and only if
  some \(x,y\in\mathcal X\) have \(\mu(x)\ne\mu(y)\) and \(r(x)\ne r(y)\);
\item if \(r\) descends, every \(x\) with \(\mu(x)=m_\ast\) has
  \(r(x)=\bar r(m_\ast)\): the selected state fixes the readout;
\item \(\mu_2\circ\tau\) descends through \(\mu\) if and only if no pair
  \(x,y\) has \(\mu(x)=\mu(y)\) and \(\mu_2(\tau x)\ne\mu_2(\tau y)\).
\end{enumerate}
\end{theorem}
\begin{proof}
(i) and (iv), with the uniqueness of \(\bar r\), are the descent
criterion of \Cref{sec:apparatus} for the surjective quotient \(\mu\),
applied to \(r\) and to \(\mu_2\circ\tau\). (ii) The first equivalence restates
the definition with the order of quantifier and negation exchanged. For the
second, if \(\bar r\) is constant then \(r=\bar r\circ\mu\) is constant; if
\(r\) is constant and \(m=\mu(x)\), \(m'=\mu(y)\), then
\(\bar r(m)=r(x)=r(y)=\bar r(m')\). Either \(\bar r\) is constant on \(M\), and
\(r\) is a law-readout, or there exist \(m,m'\) with
\(\bar r(m)\ne\bar r(m')\), and the value depends on the selected moduli
class; \(\operatorname{StateDependent}(r)\) is defined as the second case, so
exactly one holds. If \(\bar r(m)\ne\bar r(m')\), choose \(x,y\) with
\(\mu(x)=m\), \(\mu(y)=m'\); then \(r(x)\ne r(y)\) and \(\mu(x)\ne\mu(y)\).
Conversely such \(x,y\) give \(\bar r(\mu x)=r(x)\ne r(y)=\bar r(\mu y)\). (iii) If \(\mu(x)=m_\ast\), then
\(r(x)=\bar r(\mu(x))=\bar r(m_\ast)\).
\end{proof}

\subsection*{Layer instantiations}\label{layer-inst:F46}

\begingroup\small
\noindent In each item the constancy of the readout is the cited classical
result, supplied as input; \Fref{F46} then classifies that readout as a law
readout and the listed coordinates as state-dependent. It does not reprove the
cited law. In each item \(M\) is the image of the moduli quotient, so that \(\mu\) is surjective as \Cref{thm:F46} requires and the descended readout is unique.
\begin{description}
\item[\emph{Clapeyron ideal gas law (thermodynamics).}] \textit{Special case.}
For the ideal-gas model, the closure space is the
set of states \((P,V,T,n)\) with \(P,V,T,n>0\) and \(PV=nRT\), where
\(R\approx8.314\;\mathrm{J\,mol^{-1}\,K^{-1}}\) is the molar gas constant; the
moduli quotient sends a state to its \((P,V,T,n)\) coordinates; this map is
injective, so every readout descends, and the item illustrates clause (ii) only.  The readout
\(PV/(nT)\) equals \(R\) on every state of the model, so the descended readout is
constant on the moduli space and is a law-readout
in the \Fref{F46} sense.  The state-dependent readouts
(individual \(P\), \(V\), \(T\), \(n\)) vary across
moduli \citep{Clapeyron1834Memoire}.

\item[\emph{Hardy-Weinberg principle (population
genetics).}] \textit{Special case.}
For a randomly mating diploid population with two
alleles at frequencies \(p,q\) under the
Hardy-Weinberg assumptions, the closure space is
the set of equilibrium populations and the moduli
quotient collapses populations to their
allele-frequency pair \((p,q)\).  The residual readout
\((f_{AA}-p^2,\,f_{Aa}-2pq,\,f_{aa}-q^2)\) equals
\((0,0,0)\) on every equilibrium population: the
descended readout is constant and is a law-readout.
The individual genotype frequencies
\(f_{AA},f_{Aa},f_{aa}\) vary across moduli and are
state-dependent \citep{Hardy1908Mendelian}.

\item[\emph{Pythagorean theorem (Euclidean
geometry).}] \textit{Special case.}
For right triangles in Euclidean \(\mathbb R^2\),
the moduli quotient sends a triangle to its leg
lengths and hypotenuse \((a,b,c)\).  The Pythagorean
readout \(c^2-a^2-b^2\) equals \(0\) on every right
triangle; the descended readout is constant and is
a law-readout.  Triangle area \(ab/2\), the two acute
angles, and perimeter \(a+b+c\) are state-dependent
\citep{Euclid300BCElements}.

\item[\emph{Coulomb's law (electromagnetism).}] \textit{Special case.}
For a system of nonzero point charges in vacuum, the
moduli quotient sends configurations to charge values and pairwise
distances.  For each pair, with \(F\) the signed pairwise force
between charges 1 and 2 (positive for repulsion), the readout
\(F r_{12}^2/(q_1 q_2)\) equals
\(k_e=1/(4\pi\varepsilon_0)\approx8.988\times 10^9\;\mathrm{N\,m^2\,C^{-2}}\)
on every configuration of the electrostatic model of nonzero point charges at
rest in vacuum; the
descended readout is constant and is a law-readout.
Individual force magnitudes, charge values, and
distances are state-dependent
\citep{Coulomb1785Memoires}.

\item[\emph{Bayes' rule for probability updating
(probability theory).}] \textit{Special case.}
Let the closure space be the set of quadruples
\((P(H),P(D\mid H),P(D),P(H\mid D))\in[0,1]^4\) with \(P(H)>0\) and \(P(D)>0\) arising from
events \(H,D\) of some probability space, and let the moduli quotient send a
quadruple to its first three coordinates (prior, likelihood, data).  The Bayes residual
readout \(P(H\mid D)\,P(D)-P(D\mid H)\,P(H)\)
equals zero on every valid tuple; the descended
readout is constant and is a law-readout.  The
specific posterior, prior, and data values are
state-dependent \citep{Bayes1763Essay}.
\end{description}
\endgroup

\subsection{Fine-Tuning as Small Selector Region [\texorpdfstring{\Fref{F47}}{F47}]}\label{sec:self-reference-limits:fine-tuning}

The canonical physical analogy for a small target-compatible region
inside a moduli space is Weinberg's cosmological-constant analysis
\citep{Weinberg1989CosmologicalConstant}.  The layer-agnostic
statement below abstracts the smallness test from that substrate
discussion.

\paragraph{Reading the statement.} The symbols are those of \Cref{thm:F47} below.
\(\operatorname{TargetImpossible}(\Sigma(r,T))\) says that the selector
region is empty: no modulus reaches the target region. The declarations
of moduli space, target, measure and threshold are part of the setting and
add no hypothesis.

\begin{theorem}[Fine-Tuning as Small Selector Region]\label{thm:F47}
Let \(M\) be a moduli space, \(r:M\to V\) a readout, and
\(T:V\to\mathrm{Prop}\) a target-region predicate.  Let
\(\mathcal M\) be the carrier type of measure-record values, let
\(\mu:\mathrm{Pow}(M)\to\mathcal M\) be a measure record on subsets
of \(M\), let \(\mu_{\mathrm{tot}}\in\mathcal M\) be the total
measure, let \(\theta\in\Theta\) be a threshold in a declared set \(\Theta\), and let
\(\odot:\Theta\times\mathcal M\to\mathcal M\) be the declared
threshold scaling.  Let \(\operatorname{Pos}:\mathcal M\to
\mathrm{Prop}\) and \({\le_{\mathcal M}}\subseteq
\mathcal M\times\mathcal M\) be the declared positivity and order
predicates on the measure record.  Let \(\mathrm{sel}\in M\) be a
selected modulus, let \(\mathsf{selDec}\) be a selector-declared
predicate.  Define the selector region
\(\Sigma(r,T):=\{m\in M:T(r(m))\}\) and
\begin{align*}
  \operatorname{Small}(\Sigma(r,T))
  &:\Longleftrightarrow
  \operatorname{Pos}(\mu(\Sigma(r,T)))
  \wedge
  \mu(\Sigma(r,T))\le_{\mathcal M}\theta\odot\mu_{\mathrm{tot}},\\
  \operatorname{Realized}(\Sigma(r,T),\mathrm{sel},\mathsf{selDec})
  &:\Longleftrightarrow
  \mathsf{selDec}
  \wedge
  \mathrm{sel}\in\Sigma(r,T)
  \wedge
  \operatorname{Small}(\Sigma(r,T)),\\
  \operatorname{TargetImpossible}(\Sigma(r,T))
  &:\Longleftrightarrow
  \Sigma(r,T)=\emptyset.
\end{align*}
Then \(m\in\Sigma(r,T)\) if and only if \(T(r(m))\), and
\[
  \operatorname{TargetImpossible}(\Sigma(r,T))
  \Longrightarrow
  \neg\operatorname{Realized}(\Sigma(r,T),\mathrm{sel},\mathsf{selDec}):
\]
fine-tuning cannot be realized in a selector region that no modulus
reaches, whatever the measure record says; this uses only the conjunct
\(\mathrm{sel}\in\Sigma(r,T)\) of \(\operatorname{Realized}\). Moreover, realized fine-tuning
gives \(\operatorname{Pos}(\mu(\Sigma(r,T)))\) and \(\Sigma(r,T)\ne\emptyset\);
and if \(\mu\) is monotone (\(\Sigma\subseteq\Sigma'\) implies
\(\mu(\Sigma)\le_{\mathcal M}\mu(\Sigma')\)), \(\le_{\mathcal M}\) is transitive
and \(T\) implies \(T'\), then
\(\operatorname{Pos}(\mu(\Sigma(r,T)))\) and \(\operatorname{Small}(\Sigma(r,T'))\)
give \(\operatorname{Small}(\Sigma(r,T))\): a stricter target keeps a small
selector region small. If moreover \(\mathcal M=\mathbb R\) with its usual order,
\(\mu(M)=\mu(\Sigma(r,T))+\mu(\Sigma(r,\neg T))=\mu_{\mathrm{tot}}>0\),
\(\operatorname{Pos}(x):\Leftrightarrow x>0\), \(\Theta\subseteq\mathbb R\), \(\odot\) is multiplication and
\(\theta<\tfrac12\), then \(\operatorname{Small}(\Sigma(r,T))\) and
\(\operatorname{Small}(\Sigma(r,\neg T))\) do not both hold.
\end{theorem}
\begin{proof}
The membership statement is the definition of
\(\Sigma(r,T)\) as the preimage \(r^{-1}(T)\).  If
\(\operatorname{TargetImpossible}(\Sigma(r,T))\), then no \(m\in M\) lies in
\(\Sigma(r,T)\); in particular \(\mathrm{sel}\notin\Sigma(r,T)\), which
contradicts the second conjunct of \(\operatorname{Realized}\).  Hence
\(\neg\operatorname{Realized}(\Sigma(r,T),\mathrm{sel},\mathsf{selDec})\).
Realized fine-tuning contains \(\operatorname{Small}(\Sigma(r,T))\), hence
\(\operatorname{Pos}(\mu(\Sigma(r,T)))\), and \(\mathrm{sel}\in\Sigma(r,T)\).
If \(T\) implies \(T'\), then \(\Sigma(r,T)\subseteq\Sigma(r,T')\), so
monotonicity gives \(\mu(\Sigma(r,T))\le_{\mathcal M}\mu(\Sigma(r,T'))\), and
\(\operatorname{Small}(\Sigma(r,T'))\) gives
\(\mu(\Sigma(r,T'))\le_{\mathcal M}\theta\odot\mu_{\mathrm{tot}}\);
transitivity and the assumed positivity give
\(\operatorname{Small}(\Sigma(r,T))\).
For the quantitative clause, the two regions partition \(M\), so both being
small would give \(\mu_{\mathrm{tot}}\le2\theta\mu_{\mathrm{tot}}<\mu_{\mathrm{tot}}\).
\end{proof}

\subsection*{Layer instantiations}\label{layer-inst:F47}

\begingroup\small
\begin{description}
\item[\emph{Circumstellar habitable zone (astrobiology).}] \textit{Analogy.}
The moduli space is the set of orbital and
atmospheric parameter values for a planet (semi-major
axis, eccentricity, atmospheric composition, greenhouse
mass column); the readout is the equilibrium surface
temperature; the target predicate restricts to the
liquid-water range.  The selector region is the circumstellar habitable zone.
Whether it occupies a small fraction of the parameter space
depends on a declared measure and threshold satisfying
\(\operatorname{Small}\). Earth meets the liquid-water target
criterion; realizing the fine-tuning predicate additionally
requires those measure and threshold conditions
\citep{Kasting1993Habitable}.

\end{description}
\endgroup


\subsection{Topology as Global Gluing Invariant [\texorpdfstring{\Fref{F48}}{F48}]}\label{sec:self-reference-limits:topology}

The motivating mathematical example is sheaf theory, in which
topological information emerges as the consistent global gluing of
locally defined data; the standard reference is
\citet{MacLaneMoerdijk1992Sheaves}.

\begin{theorem}[Topology as Global Gluing Invariant]\label{thm:F48}
Let \(G\) be a set of global states with gluing quotient
\(\gamma:G\to\Gamma\), readout \(r:G\to V\), transport
\(\tau:G\to G'\), and target gluing quotient \(\gamma':G'\to\Gamma'\).
Assume \(\gamma\) is surjective. Write
\(\operatorname{TopologicalReadout}(r)\) when \(r\) factors through \(\gamma\),
that is \(r=\bar r\circ\gamma\) for some \(\bar r:\Gamma\to V\), and define
\[
  \mathcal O_\gamma(r)=
  \{(g,g'):\gamma(g)=\gamma(g')\ \text{and}\ r(g)\ne r(g')\},
\]
and define
\begin{align*}
  \operatorname{CompleteInvariant}(r)
  &:\Longleftrightarrow \operatorname{TopologicalReadout}(r)\wedge
    \forall g,g',\ r(g)=r(g')\Rightarrow \gamma(g)=\gamma(g'),\\
  \operatorname{CoarseInvariant}(r)
  &:\Longleftrightarrow \operatorname{TopologicalReadout}(r)\wedge
    \exists g,g',\gamma(g)\ne\gamma(g')\wedge r(g)=r(g'),\\
  \operatorname{TransportDescends}(\tau)
  &:\Longleftrightarrow
    \exists\,\bar\tau:\Gamma\to\Gamma'\ \text{ with }\
    \bar\tau\circ\gamma=\gamma'\circ\tau .
\end{align*}
Then
\[
  \operatorname{TopologicalReadout}(r)\iff\mathcal O_\gamma(r)=\emptyset,
\]
every topological readout is exactly one of a complete and a coarse
invariant, and \(\operatorname{TransportDescends}(\tau)\) holds if and only if
\(\mathcal O_\gamma(\gamma'\circ\tau)=\emptyset\).
\end{theorem}
\begin{proof}
Since \(\gamma\) is surjective, the first equivalence is the descent
criterion of \Cref{sec:apparatus}.  If \(r\) also separates \(\Gamma\)-classes, it is complete; if two
distinct gluing classes have the same \(r\)-value, it is coarse. These
are the two cases
\[
  \forall g,g',\ r(g)=r(g')\Rightarrow\gamma(g)=\gamma(g')
  \quad\text{or}\quad
  \exists g,g',\gamma(g)\ne\gamma(g')\wedge r(g)=r(g'),
\]
and each is the negation of the other, so a topological readout satisfies
exactly one of them. The transport clause is the same descent criterion applied to
\(\gamma'\circ\tau\).
\end{proof}

\subsection*{Layer instantiations}\label{layer-inst:F48}

\begingroup\small
\begin{description}
\item[\emph{Euler characteristic (algebraic
topology).}] \textit{Special case.}
For a finite cell complex, the gluing quotient
collapses cell-complex presentations to their
underlying topological space.  The Euler
characteristic
\(\chi=V-E+F\) (alternating sum of cell counts in
higher dimensions) is invariant under
cell-decomposition refinement and homotopy, hence a
TopologicalReadout.  It is a coarse invariant: the
torus and the Klein bottle both have \(\chi=0\) but
the torus is orientable and the Klein bottle is
not.  For the transport clause take \(\tau\) the identity of presentations
and \(\gamma'\) the map to homotopy types; homeomorphic spaces are homotopy
equivalent, so \(\operatorname{TransportDescends}(\tau)\) holds
\citep{Euler1750Polyhedron}.

\item[\emph{Genus of a closed orientable surface
(topology).}] \textit{Special case.}
For closed connected orientable 2-manifolds, the gluing
quotient collapses polygon-with-edge-identification
presentations to homeomorphism classes.  The genus
\(g(\Sigma)\) — the number of handles, equivalently
\((2-\chi)/2\) — is a TopologicalReadout \citep{Riemann1857Theorie}, and the later
classification theorem of closed surfaces makes
genus a CompleteInvariant for closed connected
orientable surfaces.

\item[\emph{Fundamental group (algebraic topology).}] \textit{Special case.}
For based spaces drawn from a fixed set (for instance, finite based CW
complexes), the gluing quotient sends a space to its based homotopy type.
The isomorphism class of the fundamental group \(\pi_1(X,x_0)\) — homotopy
classes of based loops — is homotopy-invariant,
hence a TopologicalReadout.  It is coarse: the
2-sphere and 3-sphere both have trivial \(\pi_1\)
but are not homotopy-equivalent.  Continuous
basepoint-preserving maps induce \(\pi_1\)-compatible
group homomorphisms \citep{Poincare1895Analysis}.

\item[\emph{{DNA} linking number of closed circular
{DNA} (molecular biology).}] \textit{Special case.}
Model a closed circular DNA duplex as an ordered pair of disjoint smooth oriented embedded circles in \(\mathbb R^3\), its two strands; the gluing quotient collapses pairs related by smooth ambient isotopy of \(\mathbb R^3\) preserving the strand orientations (no strand breaks).  The linking
number \(\mathrm{Lk}\in\mathbb Z\) — the signed
Gauss linking integral counting strand wrappings
— is invariant under smooth ambient isotopy preserving the strand orientations, hence a
TopologicalReadout.  For a duplex represented by the centreline and offset edge of an embedded closed ribbon, with compatible orientations, the Călugăreanu-White formula
\(\mathrm{Lk}=\mathrm{Tw}+\mathrm{Wr}\) decomposes
\(\mathrm{Lk}\) into non-topological twist and
writhe; only their sum is topological.
Topoisomerase enzymes implement the transport that
changes \(\mathrm{Lk}\) by strand cleavage and
religation \citep{BatesMaxwell2005DNATopology}.

\item[\emph{Reeb graph from {M}orse-function level
sets (geometric topology).}] \textit{Analogy.}
For a smooth manifold with a Morse function \(f\),
the gluing quotient collapses configurations
continuously deformable along the Morse flow.  The
Reeb graph quotients the manifold by the
``connected component of \(f^{-1}(c)\)'' relation,
encoding topology changes of level sets at
critical values.  Within the declared
function-isotopy class, the Reeb graph is a
TopologicalReadout, and a coarse invariant
(distinct Morse functions can yield isomorphic
Reeb graphs) \citep{Reeb1946Points}.
\end{description}
\endgroup

\subsection{Common-Source / Nonlocal Correlation Normal Form [\texorpdfstring{\Fref{F49}}{F49}]}\label{sec:self-reference-limits:common-source}

The Einstein--Podolsky--Rosen / Bell situation \citep{Bell1964EPR}, of correlations that admit no local common-source explanation, is an analogy, not a special case: Bell's theorem concerns probabilistic local models, whereas \(\operatorname{CommonSource}\) below is a deterministic factorization; the layer-agnostic statement below
records the descent-to-common-source obstruction abstracted from
that substrate.

\paragraph{Reading the statement.} The symbols are those of \Cref{thm:F49} below.
The locality package \(\operatorname{Loc}\) is declared data: the
theorem uses it only through its failure, which is what makes a
correlation residual nonlocal. \(\operatorname{Corr}\), by contrast, is defined from the
data: the joint readout is not a function of the two local access classes. Surjectivity of \((\sigma,q_A,q_B)\), which the theorem does not assume but a model may impose, says that every source class occurs with every pair of local access classes; it is the deterministic
form of measurement independence, and it excludes
the trivial source: for injective \(\sigma\) (for example
\(\sigma=\mathrm{id}_H\)) every fiber of \((\sigma,q_A)\) is a point, so
\(\operatorname{CommonSource}(\sigma)\) holds automatically, and surjectivity
then forces \(L_A\) and \(L_B\) to be singletons; hence when \(L_A\) or \(L_B\)
has two elements, \(\sigma\) is not injective. The
theorem tests the declared \(\sigma\), \(\iota\) and \(\rho\); it does not
assert that no source quotient explains the correlation.

\begin{theorem}[Common-Source / Nonlocal Correlation Normal Form]\label{thm:F49}
Let \(H\) be a nonempty history space with source quotient
\(\sigma:H\to S\), shared interface \(\iota:H\to I\), local access
quotients
\(q_A:H\to L_A\), \(q_B:H\to L_B\), and readouts
\(r_A:H\to A\), \(r_B:H\to B\), with joint readout
\(r=(r_A,r_B):H\to A\times B\). Let a declared direct route consist of
two declared propositions, its admissibility record and its stability record, and a route quotient
\(\rho:H\to P\). Let \(\operatorname{Corr}:\Longleftrightarrow\exists h,h'\) with
\((q_A,q_B)(h)=(q_A,q_B)(h')\) and \(r(h)\ne r(h')\), that is, the joint
readout is not fixed by the two local access classes, and let
\(\operatorname{Loc}\) be the declared locality package, a proposition the
theorem does not analyse. \(\operatorname{CommonSource}(\sigma)\) below is a
deterministic factorization; probabilistic common-cause models are not
covered. No surjectivity is assumed. Write \(\operatorname{CommonSource}(\sigma)\) when \(r_A\) is
constant on the fibers of \((\sigma,q_A)\) and \(r_B\) is constant on the
fibers of \((\sigma,q_B)\); \(\operatorname{InterfaceFactorization}(\iota)\)
when \(r\) is constant on the fibers of \((\iota,q_A,q_B)\); and
\(\operatorname{DirectRoute}\) when both records hold and
\(r\) is constant on the fibers of \(\rho\), and define
\[
  \operatorname{LocallyExplainable}
  :\iff \operatorname{CommonSource}(\sigma)\vee\operatorname{DirectRoute}
    \vee \operatorname{InterfaceFactorization}(\iota).
\]
Then
\begin{align*}
  \operatorname{CommonSource}(\sigma)
  &\iff
  \exists\,\bar r_A:S\times L_A\to A,\ \bar r_B:S\times L_B\to B,\\
  &\hphantom{{}\iff{}}
    \bar r_A\circ(\sigma,q_A)=r_A\ \wedge\ \bar r_B\circ(\sigma,q_B)=r_B,\\
  \operatorname{InterfaceFactorization}(\iota)
  &\iff
  \exists\,\hat r:I\times L_A\times L_B\to A\times B,\
    \hat r\circ(\iota,q_A,q_B)=r,\\
  \operatorname{DirectRoute}
  &\iff
  \text{both records hold, and}\
    \exists\,\tilde r:P\to A\times B,\ \tilde r\circ\rho=r.
\end{align*}
Call \(\operatorname{Corr}\wedge\neg\operatorname{LocallyExplainable}\) a
\emph{correlation residual}, and a \emph{nonlocal residual} when moreover
\(\neg\operatorname{Loc}\). By the three equivalences, a correlation residual
is present exactly when \(\operatorname{Corr}\) holds, neither the common-source
nor the interface factorization exists, and it is not the case that the route
is admissible, stable and has some \(\tilde r\) with \(\tilde r\circ\rho=r\).
Moreover, \(\neg\operatorname{Corr}\) implies \(\operatorname{InterfaceFactorization}(\iota)\)
for every \(\iota\), so a correlation residual is exactly
\(\neg\operatorname{LocallyExplainable}\); and \(\operatorname{CommonSource}(\sigma)\)
implies \(\operatorname{InterfaceFactorization}\) for the interface \(\iota:=\sigma\).
\end{theorem}
\begin{proof}
Since \(H\) is nonempty, so are \(A\) and \(B\); a map on \(H\) that is constant on the fibers of a map \(p\) factors through \(p\), defined from representatives on the image of \(p\) and by a fixed value elsewhere. (Without nonemptiness the CommonSource equivalence fails: for \(H=L_B=A=\varnothing\) and \(S\times L_A\) a point, \(\operatorname{CommonSource}(\sigma)\) holds vacuously but no \(\bar r_A\) exists.) So a map on \(H\) with values in \(A\), \(B\) or \(A\times B\) factors through a map \(p\) exactly when it is constant on the fibers of \(p\). Hence \(r_A\) and \(r_B\) are constant on the
fibers of \((\sigma,q_A)\) and \((\sigma,q_B)\) exactly when response maps
\(\bar r_A\), \(\bar r_B\) as displayed exist: each local readout depends only
on the shared source class and its own access data. The
joint readout factors through \((\iota,q_A,q_B)\)
exactly when it is constant on its fibers, which gives the interface
case. The direct-route case is the same factorization argument for \(\rho\): the joint readout factors through the route
quotient exactly when it is constant on its fibers, and a declared route
explains the correlation when, in addition, its admissibility and
stability records hold. A correlation is
locally explained exactly when one of the three explanations is
supplied. If a correlation is present and none is supplied, it is
recorded as a residual, not as an explained relation; it is a nonlocal
residual only relative to a failed locality package.
If \(\neg\operatorname{Corr}\), then \(r\) is constant on the fibers of \((q_A,q_B)\), hence on the finer fibers of \((\iota,q_A,q_B)\). Under \(\operatorname{CommonSource}(\sigma)\), \(r_A\) and \(r_B\) are constant on the fibers of \((\sigma,q_A)\) and \((\sigma,q_B)\), so \(r=(r_A,r_B)\) is constant on the fibers of \((\sigma,q_A,q_B)\).
\end{proof}

\subsection*{Layer instantiations}\label{layer-inst:F49}

\begingroup\small
\begin{description}
\item[\emph{{MISTRA} twin studies (behavioral
genetics).}] \textit{Analogy.}
For the Minnesota Study of Twins Reared Apart, the
source quotient sends a twin pair to its
shared-genotype class; the readouts are conditional
trait
distributions
\(P(\mathrm{trait}\mid\sigma)\) under a declared
quantitative-genetic model.  The CommonSource
hypothesis holds when conditional distributions
descend through the shared-genotype quotient; the
interface captures shared rearing environment
(absent for the reared-apart sample).  The reared-apart monozygotic
correlation \(r_{MZA}\), which estimates heritability directly,
operationalises the \Fref{F49} LocallyExplainable decomposition; the
Falconer formula \(h^2=2(r_{MZ}-r_{DZ})\) is the counterpart for
twins reared together
\citep{Bouchard1990Sources}.

\end{description}
\endgroup


\subsection{Vacuum / Background Selection [\texorpdfstring{\Fref{F50}}{F50}]}\label{sec:self-reference-limits:vacuum}

The most discussed contemporary substrate analogy is the string
landscape of vacua surveyed in \citet{Susskind2003Landscape};
the layer-agnostic statement below records the selector-singleton
shape without committing to that physical realization.

\begin{theorem}[Vacuum / Background Selection]\label{thm:F50}
Let \(\mathcal X\) be a closure space with moduli quotient
\(\mu:\mathcal X\to M\), raw background \(b_0:\mathcal X\to B\),
a declared background \(b:M\to B\), target \(t:M\to T\), and vacuum
criterion \(V\subseteq M\), and a declared set \(T_{\mathrm{f}}\subseteq T\)
of formed targets. Assume \(\mu\) is surjective. Define
\begin{align*}
  \operatorname{NecessaryBackground}(b)
  &:\Longleftrightarrow \forall m,m'\in M,\ b(m)=b(m'),\\
  \operatorname{BackgroundModuli}(b)
  &:\Longleftrightarrow \exists m,m'\in M,\ b(m)\ne b(m'),\\
  \operatorname{Vacuum}(m_\ast)
  &:\Longleftrightarrow V=\{m_\ast\}\ \wedge\ t(m_\ast)\in T_{\mathrm{f}}.
\end{align*}
Then
\[
  \bigl(\forall x,y\in\mathcal X,\ \mu(x)=\mu(y)\Rightarrow b_0(x)=b_0(y)\bigr)
  \iff
  \exists\,\beta:M\to B\ \text{ with } \beta\circ\mu=b_0.
\]
The necessary/moduli alternatives are exhaustive and disjoint. If \(b\) is
a descended raw background, \(b\circ\mu=b_0\), then
\[
  \operatorname{NecessaryBackground}(b)
  \iff \forall x,y\in\mathcal X,\ b_0(x)=b_0(y).
\]
Exactly one of \(V=\emptyset\), \(V\) a singleton, and \(V\) containing two
distinct elements holds. Some \(m_\ast\) satisfies
\(\operatorname{Vacuum}(m_\ast)\) if and only if \(V=\{m\}\) for some \(m\)
with \(t(m)\in T_{\mathrm{f}}\), and such an \(m_\ast\) is unique. If
\(\operatorname{Vacuum}(m_\ast)\), then under \(\operatorname{BackgroundModuli}(b)\) some modulus outside \(V\) carries a background different from \(b(m_\ast)\); if moreover \(b\circ\mu=b_0\), then \(b_0(x)=b(m_\ast)\) for
every \(x\) with \(\mu(x)=m_\ast\).
\end{theorem}
\begin{proof}
Since \(\mu\) is surjective, the descent clause is the descent criterion
of \Cref{sec:apparatus} for \(\mu\) and \(b_0\).  For the declared background \(b\), either \(b\) is constant on \(M\)
or there exist \(m,m'\) with \(b(m)\ne b(m')\):
\[
  \bigl(\forall m,m', b(m)=b(m')\bigr)
  \quad\vee\quad
  \bigl(\exists m,m', b(m)\ne b(m')\bigr).
\]
These alternatives are mutually exclusive. For the vacuum criterion,
either \(V\) is empty, or it has an element \(m\); in the second case either
every element of \(V\) equals \(m\), so \(V=\{m\}\), or some element differs
from \(m\), so \(V\) contains two distinct elements. A singleton has no two
distinct elements and a nonempty set is not empty, so exactly one case holds.
By definition, \(\operatorname{Vacuum}(m_\ast)\) says that \(V=\{m_\ast\}\) and
\(t(m_\ast)\in T_{\mathrm{f}}\), which gives the existence equivalence; if
\(\operatorname{Vacuum}(m_\ast)\) and \(\operatorname{Vacuum}(m'_\ast)\), then
\(\{m_\ast\}=V=\{m'_\ast\}\), so \(m_\ast=m'_\ast\).
If \(b\circ\mu=b_0\), then \(b_0(x)=b(\mu x)\), so \(b_0\) is constant when
\(b\) is; conversely every \(m\in M\) is \(\mu(x)\) for some \(x\), so \(b\) is
constant when \(b_0\) is. For the last clause, \(b_0(x)=b(\mu x)=b(m_\ast)\) when
\(\mu(x)=m_\ast\); and if \(b(m)\ne b(m')\), one of \(m,m'\) has background
different from \(b(m_\ast)\), and it is not \(m_\ast\), the only member of \(V\).
\end{proof}

\subsection*{Layer instantiations}\label{layer-inst:F50}

\begingroup\small
\begin{description}
\item[\emph{Waddington epigenetic landscape
(developmental biology).}] \textit{Analogy.}
For a developing embryo, the moduli quotient
collapses stem-cell trajectory configurations
sharing the same cell-type endpoint; the
descended background records the cell-fate class.
Waddington's epigenetic landscape depicts cell-fate
decision as a ball rolling down a branching
landscape: each branching point is a vacuum-style
selection of one cell-fate from a set of accessible
attractors \citep{Waddington1957Strategy}.

\end{description}
\endgroup


\subsection{Unification as Common Refinement [\texorpdfstring{\Fref{F51}}{F51}]}\label{sec:self-reference-limits:unification}

The motivating particle-physics analogy is the electroweak unification
introduced in \citet{Weinberg1967Leptons}.  The statement below records a common-refinement shape; electroweak unification is an analogy, and no physical unification is asserted to satisfy its hypotheses.

\paragraph{Reading the statement.}
\emph{Exact unification} is defined as the conjunction displayed in the
theorem: formation of the parent closure, the three pairs of equations, role
preservation of the two child theories, and the status check of every cross
residual. Role preservation and the residual status check are declared
predicates of the apparatus. The content of the theorem is what follows from
the first pair of equations: the parent quotient is a common refinement of the
two children, in the sense of the displayed square, and, when \(q_U\) is
surjective, \(\langle\pi_A,\pi_B\rangle\) is injective exactly when the parent
carries no distinction beyond those of the two children.

\begin{theorem}[Unification as Common Refinement]\label{thm:F51}
Let \(\mathcal C_U,\mathcal C_A,\mathcal C_B\) be parent and child
closures with quotients \(q_U:\mathcal C_U\to Q_U\),
\(q_A:\mathcal C_A\to Q_A\), and \(q_B:\mathcal C_B\to Q_B\). Let
\(\iota_A:\mathcal C_U\to\mathcal C_A\) and
\(\iota_B:\mathcal C_U\to\mathcal C_B\) be interpretations, and let
\(\pi_A:Q_U\to Q_A\), \(\pi_B:Q_U\to Q_B\) be quotient projections.
Let \(R_U^A,R_U^B:Q_U\to Q_U\) and \(R_A:Q_A\to Q_A\),
\(R_B:Q_B\to Q_B\) be the declared route systems. For claim sets
\(K_A,K_B,K_U\), lifts \(\lambda_A:K_A\to K_U\),
\(\lambda_B:K_B\to K_U\), status maps \(s_A,s_B,s_U\), and status
translations \(\alpha_A,\alpha_B\) into the statuses of \(K_U\), \emph{exact
unification} is defined as the conjunction of the following: the
declared formedness proposition for the parent closure \(\mathcal C_U\) holds; the equations
\begin{align*}
  \pi_A\circ q_U=q_A\circ\iota_A
  &\quad\wedge\quad
  \pi_B\circ q_U=q_B\circ\iota_B,\\
  \pi_A\circ R_U^A=R_A\circ\pi_A
  &\quad\wedge\quad
  \pi_B\circ R_U^B=R_B\circ\pi_B,\\
  s_U(\lambda_A k)=\alpha_A(s_A k)
  &\quad\wedge\quad
  s_U(\lambda_B \ell)=\alpha_B(s_B \ell)
  \quad\text{for all }k\in K_A,\ell\in K_B
\end{align*}
hold; the declared role-preservation proposition holds; and the declared cross-residual status-check proposition holds. If
\(\pi_A\circ q_U=q_A\circ\iota_A\) and \(\pi_B\circ q_U=q_B\circ\iota_B\), in
particular under exact unification, then
\[
  \langle\pi_A,\pi_B\rangle\circ q_U
  =(q_A\times q_B)\circ\langle\iota_A,\iota_B\rangle
  :\mathcal C_U\to Q_A\times Q_B ,
\]
and the parent kernel is contained in both pulled-back child kernels: for all
\(x,x'\in\mathcal C_U\), \(q_U(x)=q_U(x')\) implies
\(q_A(\iota_A x)=q_A(\iota_A x')\) and \(q_B(\iota_B x)=q_B(\iota_B x')\).
If moreover \(q_U\) is surjective, then \(\langle\pi_A,\pi_B\rangle\) is injective
if and only if \(q_U(x)=q_U(x')\Leftrightarrow q_A(\iota_A x)=q_A(\iota_A x')
\wedge q_B(\iota_B x)=q_B(\iota_B x')\) for all \(x,x'\): the parent carries
exactly the distinctions of the two children.
\end{theorem}
\begin{proof}
The first pair of equations gives, for every
\(x\in\mathcal C_U\),
\[
  \pi_A(q_U x)=q_A(\iota_A x),
  \qquad
  \pi_B(q_U x)=q_B(\iota_B x),
\]
which is the displayed square componentwise, so \(Q_U\) is a common
refinement. If \(q_U(x)=q_U(x')\), then
\(q_A(\iota_A x)=\pi_A(q_U x)=\pi_A(q_U x')=q_A(\iota_A x')\), and likewise
for \(B\). For the last clause, let \(q_U\) be surjective. By the square,
\(q_A(\iota_A x)=q_A(\iota_A x')\wedge q_B(\iota_B x)=q_B(\iota_B x')\) says
\(\langle\pi_A,\pi_B\rangle(q_U x)=\langle\pi_A,\pi_B\rangle(q_U x')\). If the
pairing is injective, this gives \(q_U(x)=q_U(x')\). Conversely, every two
points of \(Q_U\) are \(q_U(x)\) and \(q_U(x')\) for some \(x,x'\), so the
equivalence gives injectivity.
\end{proof}

\subsection*{Layer instantiations}\label{layer-inst:F51}

\begingroup\small
\begin{description}
\item[\emph{Maxwell unification of electromagnetism
(classical physics).}] \textit{Analogy.}
The parent closure is Maxwell's electromagnetic
field theory; the children are pre-Maxwellian
electrostatics (Coulomb-Gauss) and magnetostatics
(Amp\`ere-Biot-Savart).  Setting time derivatives to
zero in Maxwell's equations yields each child as a
specialisation.  The unified parent predicts
electromagnetic waves at the speed of light,
identifying light as electromagnetic — a
prediction not derivable from either child alone
\citep{Maxwell1865Dynamical}.

\end{description}
\endgroup


\subsection{Closure Completeness Boundary [\texorpdfstring{\Fref{F52}}{F52}]}\label{sec:self-reference-limits:closure-completeness}

A proof-theoretic analogue of a closure-completeness boundary is
G\"odel's incompleteness theorem \citep{Godel1931Unentscheidbare}, where
incompleteness is derived from the arithmetic of the theory; the statement
below records a declared scope boundary instead. The modal-logic systematization of provability obstructions is
collected in \citet{Boolos1993LogicProvability}.  The layer-agnostic
meta-coherence statement below records the apparatus-side scope
discipline that complements those classical results.

\begin{definition}[Catalog boundary]\label{def:catalog-boundary}
Let \(D\) be a direction type with predicates \(\operatorname{claimed}\),
\(\operatorname{inScope}\) and \(\operatorname{namedOutside}\), as in
\Cref{thm:F52}. The \emph{catalog boundary} is
\(\partial D:=\{d:\operatorname{claimed}(d)\wedge\neg\operatorname{inScope}(d)
\wedge\operatorname{namedOutside}(d)\}\), the claimed directions explicitly
marked out of scope, for example because they require substrate-specific data
or stronger apparatus. The \emph{unmarked frontier} is
\(\{d:\operatorname{claimed}(d)\wedge\neg\operatorname{inScope}(d)\wedge
\neg\operatorname{namedOutside}(d)\}\).
\end{definition}

The theorem below states when a closure package reaches this boundary:
it is complete up to its declared scope, and every direction it claims
lies inside that scope or is named as outside it.

\paragraph{Reading the statement.}
\(\operatorname{formed}\) records that the closure package is formed in
the sense of \Cref{def:closure-formation}, and
\(\operatorname{scopeDeclared}\) records that its completeness scope has
been declared. Both are declared data of the package. Each displayed
line of the theorem unfolds the predicate on its left; together they
state what completeness up to a declared scope requires. The last sentence of the theorem, that each conjunct can fail while the other three hold, is proved by one-element examples.

\begin{theorem}[Closure Completeness Boundary]\label{thm:F52}
Let \(R\) be a residual type with predicates
\(\operatorname{hasStatus}:R\to\mathrm{Prop}\) and
\(\operatorname{adequate}:R\to\mathrm{Prop}\). Let \(D\) be a direction
type with predicates \(\operatorname{inScope}:D\to\mathrm{Prop}\),
\(\operatorname{claimed}:D\to\mathrm{Prop}\), and
\(\operatorname{namedOutside}:D\to\mathrm{Prop}\). Let
\(\operatorname{formed}\) and \(\operatorname{scopeDeclared}\) be
propositions, and let
\(v\) be a meta-audit verdict, drawn from a declared set of verdicts
containing \(\operatorname{accepted}\). Define
\begin{align*}
  \operatorname{ResidualCoverage}
  &:\Longleftrightarrow
  \forall r\in R,\ \operatorname{hasStatus}(r)\wedge\operatorname{adequate}(r),\\
  \operatorname{MetaCovered}(v)
  &:\Longleftrightarrow v=\operatorname{accepted},\\
  \operatorname{ScopedClosureComplete}
  &:\Longleftrightarrow
  \operatorname{formed}\wedge\operatorname{scopeDeclared}\wedge
  \operatorname{ResidualCoverage}\\
  &\qquad\wedge\
  \forall d\in D,\ \operatorname{claimed}(d)\Rightarrow
    \bigl(\operatorname{inScope}(d)\vee\operatorname{namedOutside}(d)\bigr),\\
  \operatorname{MetaComplete}
  &:\Longleftrightarrow \operatorname{ScopedClosureComplete}\wedge\operatorname{MetaCovered}(v).
\end{align*}
Then \(\operatorname{ScopedClosureComplete}\) holds if and only if
\(\operatorname{formed}\) and \(\operatorname{scopeDeclared}\) hold and the
three obstruction sets
\[
\begin{gathered}
  \{r:\neg\operatorname{hasStatus}(r)\},\qquad
  \{r:\neg\operatorname{adequate}(r)\},\\
  \{d:\neg\operatorname{inScope}(d)\wedge\operatorname{claimed}(d)
      \wedge\neg\operatorname{namedOutside}(d)\}
\end{gathered}
\]
are empty; the third set is the unmarked frontier of
\Cref{def:catalog-boundary}. Hence \(\operatorname{MetaComplete}\) fails
exactly when one of these conditions fails or \(v\ne\operatorname{accepted}\).
Each of the four conjuncts of \(\operatorname{ScopedClosureComplete}\) can fail
while the other three hold.
\end{theorem}
\begin{proof}
\(\operatorname{ResidualCoverage}\) says that every residual has a status and
that every residual is adequate, that is, that the first two sets are empty.
With \(\operatorname{hasStatus}(r):\Longleftrightarrow\exists\rho\in\mathrm{Rec},\ \operatorname{loc}(\rho)=r\),
its first conjunct is (RC) of \Cref{sec:apparatus}; adequacy is an additional
requirement.
For a direction \(d\) there are three cases:
\[
  \operatorname{inScope}(d),\qquad
  \operatorname{namedOutside}(d),\qquad
  \neg\operatorname{inScope}(d)\wedge\neg\operatorname{namedOutside}(d).
\]
The direction clause of \(\operatorname{ScopedClosureComplete}\) excludes
the third case for every claimed direction, which is emptiness of the third
set. The penultimate sentence follows from the definition of
\(\operatorname{MetaComplete}\). For the last, take \(R\) and \(D\) with one
element each. Making \(\operatorname{formed}\) false, or
\(\operatorname{scopeDeclared}\) false, or \(\operatorname{hasStatus}\) false
at the residual, or \(d\) claimed with \(\operatorname{inScope}(d)\) and
\(\operatorname{namedOutside}(d)\) false, while every other datum is true,
falsifies exactly one conjunct.
\end{proof}

For the catalog itself, take as \(D\) the candidate directions
considered for it, let \(\operatorname{inScope}(d)\) mean that \(d\) is a
law of the catalog or reduces to one, let
\(\operatorname{namedOutside}(d)\) mean that \(d\) is named in
\Cref{sec:sketches}, and let \(\operatorname{claimed}(d)\) hold for
every \(d\in D\).  The direction clause then requires every candidate to
be covered or named, and the residual clauses require a residual
register for the catalog, which is not given here; \Cref{sec:sketches}
discusses this application.

\subsection*{Layer instantiations}\label{layer-inst:F52}

\begingroup\small
\begin{description}
\item[\emph{{CONSORT} statement for randomised-trial
reporting (clinical research).}] \textit{Analogy.}
The residuals are the methodological items required
for adequate trial reporting; directions are
declarable claims and outcomes.  CONSORT 2010
requires 25 explicit reporting items, with primary
and secondary outcomes \emph{inScope} and
exploratory analyses \emph{namedOutside}.  Residual
coverage holds when every CONSORT item has status
and is adequate; the meta-audit is journal-editor
acceptance \citep{Schulz2010CONSORT}.

\end{description}
\endgroup


\clearpage
\subsection*{Self-Reference and Limits at a glance}\label{sec:self-reference-limits:summary}\addcontentsline{toc}{subsection}{Self-Reference and Limits at a glance}
\begin{center}
  \centering
  \captionof{table}{Self-Reference and Limits at a glance: the eighteen laws of \Cref{sec:self-reference-limits}, grouped into reading groups.}
  \label{tab:axis-self-reference-limits-summary}
  \scriptsize
  \setlength{\tabcolsep}{4pt}
  \medskip\noindent\begin{minipage}{\textwidth}\centering
  \textbf{Limits}\par\smallskip
  \begin{tabularx}{\textwidth}{@{}
      >{\raggedright\arraybackslash}p{0.24\textwidth}
      >{\raggedright\arraybackslash}X
      >{\raggedright\arraybackslash}p{0.10\textwidth}@{}}
    \toprule
    Name & One-line claim & Substrate cite \\
    \midrule
    Horizon-Sufficiency Theorem (\Fref{F33}) &
    When the pairing is surjective, a target depends on a history only
    through its boundary readout and
    exterior class exactly when it factors through their pairing; a
    bundled family does so exactly when each member does.  &
    \textemdash \\
    Information-Loss Theorem (\Fref{F34}) &
    When recovery after transport is surjective, source information is
    recoverable after transport exactly when it
    agrees on any two sources that transport and recovery identify.  &
    \textemdash \\
    Singular-Limit Theorem (\Fref{F36}) &
    For a surjective limit quotient, a target is completed exactly when no two processes with the same
    limit object have different target values, and is path-dependent
    otherwise; the pairing \(\rho_{u,\ell}\) refines \(\ell\), and \(u\) is completed
    through it; \(R\circ u\) is completed when \(u\) is, conversely for injective
    \(R\).  &
    \textemdash \\
    Locality Theorem (\Fref{F45}) &
    For a surjective local quotient, a target is local to a domain exactly
    when it is constant on the domain's classes. For strict coarsening, with the signature in the
    declared class and a nonempty carrier, a domain is minimal exactly
    when it has the signature's fibers.  &
    \textemdash \\
    \bottomrule
  \end{tabularx}
  \end{minipage}\par

  \medskip\noindent\begin{minipage}{\textwidth}\centering
  \textbf{Emergence}\par\smallskip
  \begin{tabularx}{\textwidth}{@{}
      >{\raggedright\arraybackslash}p{0.24\textwidth}
      >{\raggedright\arraybackslash}X
      >{\raggedright\arraybackslash}p{0.10\textwidth}@{}}
    \toprule
    Name & One-line claim & Substrate cite \\
    \midrule
    Scale-Descent Theorem (\Fref{F35}) &
    For a surjective scale quotient and coarsening, a law descends to the
    coarser scale exactly when it descends to the
    finer scale and its descended law is constant on the fibers of the
    coarsening; the renormalized law is then unique.  &
    \textemdash \\
    Continuum-Emergence Theorem (\Fref{F43}) &
    Over a nonempty directed index set with a surjective limit quotient,
    a coherent, cofinal tower with vanishing residuals, compatible
    projections and stage readouts, and a target readout equal to one
    stage readout after projection and constant on the limit fibers
    yields \(\operatorname{StableLimitQuotient}(L)\), witnessed by the unique readout on the quotient agreeing with every stage readout; coherence, cofinality and vanishing enter only as conjuncts of StableLimitQuotient, and the content is descent of the common stage readout through \(\ell\).  &
    \textemdash \\
    Fine-Tuning Theorem (\Fref{F47}) &
    If no modulus reaches the target region, fine-tuning is not
    realized; for a monotone measure record with a transitive order, a
    stricter target whose selector region has positive measure keeps a small
    selector region small.  &
    \textemdash \\
    Vacuum Theorem (\Fref{F50}) &
    For a surjective moduli quotient and a descended background, exactly
    one holds: it is necessary
    (equivalently, the raw background is constant) or it varies across
    moduli; a vacuum exists exactly when the vacuum criterion is a
    singleton with formed target, and is then unique.  &
    \textemdash \\
    Unification Theorem (\Fref{F51}) &
    Under exact unification the parent quotient is a common refinement of
    both children; when \(q_U\) is surjective, \(\langle\pi_A,\pi_B\rangle\) is
    injective iff the parent keeps exactly the children's distinctions.  &
    \textemdash \\
    \bottomrule
  \end{tabularx}
  \end{minipage}\par

  \medskip\noindent\begin{minipage}{\textwidth}\centering
  \textbf{Invariants}\par\smallskip
  \begin{tabularx}{\textwidth}{@{}
      >{\raggedright\arraybackslash}p{0.24\textwidth}
      >{\raggedright\arraybackslash}X
      >{\raggedright\arraybackslash}p{0.10\textwidth}@{}}
    \toprule
    Name & One-line claim & Substrate cite \\
    \midrule
    Complementarity Theorem (\Fref{F37}) &
    On a nonempty carrier, two separately admissible accesses are
    complementary exactly when
    every admissible joint carrier identifies two histories that one
    access separates, so that no admissible carrier has both accesses
    factoring through it.  &
    \textemdash \\
    Arrow Theorem (\Fref{F38}) &
    For a surjective record quotient, a continuation has an arrow exactly
    when it descends to a map of record classes that never lowers a
    record's status.  &
    \textemdash \\
    Entropy Theorem (\Fref{F39}) &
    For a surjective access quotient, the target obstruction is empty
    exactly when the target descends. For an entropy functional with a declared zero satisfying
    \(E(q')=0\) exactly when \(q'\) is injective,
    and with the chain rule and combine-monotonicity at a coarsening
    pair, \(E(q)=0\) exactly when \(q\) has no raw split pair, and entropy
    does not decrease under the coarsening. For finite nonempty \(H\), the fiber-volume functional
    \(E(q')=|H|^{-1}\sum_h\log|q'^{-1}(q'h)|\) satisfies these hypotheses and has \(m(\langle q,t\rangle,q)=0\) exactly
    when the target descends.  &
    \textemdash \\
    Irreducibility Theorem (\Fref{F42}) &
    Closedness at \(h\) is emptiness of the target and route split-pair sets on the class of \(h\); with well-founded costs and some admissible closed compression, an irreducible one exists; if for every \(v\in E\) the admissible class contains a map taking \(h\) to \(v\) and isolating \(h\), irreducibility is minimal cost over \(E\) and ignores the target.  &
    \textemdash \\
    Criticality Theorem (\Fref{F44}) &
    A phase change along a finite path of formed controls in which, for every neighbourhood index, each control has a neighbour in the phase of the next, or along a continuous path when neighbourhoods are the open sets, crosses a phase boundary; boundary and criticality
    descend through the control quotient when their data do.  &
    \textemdash \\
    Law--State Split Theorem (\Fref{F46}) &
    For a surjective moduli quotient, a readout descends exactly when no two closures with one modulus differ, and then uniquely; a descended readout is a law readout exactly when it is constant on the closure space and state-dependent exactly when two closures with different moduli differ; the selected modulus fixes the readout on its fiber.  &
    \textemdash \\
    Topology Theorem (\Fref{F48}) &
    For a surjective gluing quotient, a readout factors through it exactly when its split-pair set is empty; each such readout is exactly one of complete and coarse; a transport induces a map of gluing classes exactly when \(\gamma'\circ\tau\) has empty split-pair set.  &
    \textemdash \\
    Common-Source Theorem (\Fref{F49}) &
    On a nonempty history space, common-source and interface
    explanations are equivalent to their factorizations; direct-route
    explanation additionally requires admissibility and stability. A
    correlation residual is exactly failure of local explainability, and a
    common source supplies an interface factorization with \(\iota=\sigma\).  &
    \textemdash \\
    Closure Completeness Theorem (\Fref{F52}) &
    Scoped closure completeness holds exactly when the closure is formed
    and scope-declared, every residual is statused and adequate, and
    every claimed direction is in scope or named outside.  &
    \textemdash \\
    \bottomrule
  \end{tabularx}
  \end{minipage}\par

\end{center}

\FloatBarrier
